\documentclass[11pt]{article}
\usepackage{amssymb}
\usepackage{amsthm}
\usepackage{enumitem}
\usepackage{amsmath}
\usepackage{bm}
\usepackage{adjustbox}
\usepackage{mathrsfs}
\usepackage{graphicx}
\usepackage{physics}
\usepackage{siunitx}
\usepackage[mathscr]{euscript}

\title{\textbf{Solved selected problems of Introductory functional analysis with applications - Erwin Kreyszig}}
\author{Franco Zacco}
\date{}

\DeclareSymbolFontAlphabet{\mathscrsfs}{rsfs}
\addtolength{\topmargin}{-3cm}
\addtolength{\textheight}{3cm}

\newcommand{\N}{\mathbb{N}}
\newcommand{\Z}{\mathbb{Z}}
\newcommand{\Q}{\mathbb{Q}}
\newcommand{\R}{\mathbb{R}}
\newcommand{\C}{\mathbb{C}}
\newcommand{\diam}{\text{diam}}
\newcommand{\cl}{\text{cl}}
\newcommand{\bdry}{\text{bdry}}
\newcommand{\inter}{\text{int}}
\newcommand{\dom}{\mathscrsfs{D}}
\newcommand{\range}{\mathscrsfs{R}}
\newcommand{\nullsp}{\mathscrsfs{N}}
\newcommand{\inprd}[1]{\left\langle{#1}\right\rangle}
\newcommand{\re}[1]{\text{Re}{#1}}
\newcommand{\im}[1]{\text{Im}{#1}}
\newcommand{\Span}[1]{\text{span}(#1)}


\theoremstyle{definition}
\newtheorem*{solution*}{Solution}

\begin{document}
\maketitle
\thispagestyle{empty}

\section*{Chapter 3 - Inner Product Spaces. Hilbert Spaces}
\subsection*{3.1 - Inner Product Space. Hilbert Space}
\begin{proof}{\textbf{1}}
Let us prove (4) as follows
\begin{align*}
    \|x+y\|^2 + \|x-y\|^2 &= \inprd{x+y, x+y} + \inprd{x-y, x-y}\\
    &=\inprd{x, x} + \inprd{x, y} + \inprd{y, x} + \inprd{y, y}
    + \inprd{x, x} - \inprd{x, y} - \inprd{y, x} + \inprd{y, y}\\
    &=2\inprd{x, x} + 2\inprd{y, y}\\
    &=2(\|x\|^2 + \|y\|^2)
\end{align*}
Where we used the definition of the norm $\|x\| = \sqrt{\inprd{x,x}}$ and
the distributive property of the inner products.
\end{proof}
\begin{proof}{\textbf{2}}
Let $x \perp y$ in an inner product space $X$ then
\begin{align*}
    \|x + y\|^2 &= \inprd{x+y, x+y}\\
    &= \inprd{x, x} + \inprd{x,y} + \inprd{y,x} + \inprd{y,y}\\
    &= \inprd{x, x} + \inprd{y, y}\\
    &= \|x\|^2 + \|y\|^2
\end{align*}
Where we used that $\inprd{x,y} = \inprd{y,x} = 0$ since $x \perp y$.

Suppose now we have $m$ mutually orthogonal vectors, i.e.
$x_1 \perp x_2 \perp ... \perp x_m$ then
\begin{align*}
    \|x_1 + x_2 + ... + x_m\|^2 &= \inprd{x_1 + x_2 + ... + x_m, x_1 + x_2 + ... + x_m}\\
    &= \inprd{x_1, x_1} + \inprd{x_1,x_2} +  ...  +\inprd{x_m,x_{m-1}} + \inprd{x_m,x_m}\\
    &= \inprd{x_1, x_1} + \inprd{x_2, x_2} + ... + \inprd{x_m, x_m}\\
    &= \|x_1\|^2 + \|x_2\|^2 + ... + \|x_m\|^2
\end{align*}
Where we used that $\inprd{x_i,x_j} = 0$ for $i \neq j$ since they are
orthogonal.
\end{proof}

\cleardoublepage
\begin{proof}{\textbf{3}}
Let $x,y \in \R$ and suppose we have the following property
$$\|x + y\|^2 = \|x\|^2 + \|y\|^2$$
then since $\|x\|^2 = \inprd{x,x}$ above equation implies that
\begin{align*}
    \inprd{x +y, x +y} &= \inprd{x,x} + \inprd{y,y}\\
    \inprd{x, x} + \inprd{x,y} + \inprd{y,x} + \inprd{y,y}
    &=\inprd{x,x} + \inprd{y,y}\\
    \inprd{x,y} + \inprd{y,x} &= 0
\end{align*}
But since $x,y \in \R$ we have that $\inprd{x,y} = \inprd{y,x}$, hence
$$\inprd{x,y} = \inprd{y,x} = 0$$
Therefore $x \perp y$.

Let now $x,y \in \C$ and suppose that the property holds in this case too.
Then as before must be that
\begin{align*}
    \inprd{x,y} + \inprd{y,x} &= 0
\end{align*}
Also, we know that $\inprd{y,x} = \overline{\inprd{x,y}}$ hence
we need that
\begin{align*}
    \inprd{x,y} &= -\overline{\inprd{x,y}}
\end{align*}
But since we are assuming a complex inner product we have that $\inprd{x,y}$
is of the form $a + bi$ hence
\begin{align*}
    a + bi = -(a - bi) = -a + bi
\end{align*}
Equating the real and imaginary parts we see that $a = 0$ and $b = b$.
So for example if $x = 1 + i$ and $y = 1 - i$ we see that the Pythagorean
property holds, as we show below
\begin{align*}
    \|1 + i + 1 -i\|^2 = \|2\|^2
    = 4 = 2 + 2 = \|1 + i\|^2 + \|1 - i\|^2 
\end{align*}
But inner product gives us
$$\inprd{x,y} = (1 + i)\overline{(1 - i)} = (1 + i)(1 + i)
= 1 + i + i -1 = 2i$$
Therefore $x$ and $y$ are not orthogonal.
\end{proof}

\cleardoublepage
\begin{proof}{\textbf{6}}
Let $x \neq 0$ and $y \neq 0$
\begin{itemize}
\item [(a)] Let $x \perp y$ we want to prove that $\{x,y\}$ is a linearly
independent set i.e. we want to prove
\begin{align*}
    \alpha_1 x + \alpha_2 y = 0
\end{align*}
Only for $\alpha_1 = \alpha_2 = 0$.
Suppose the opposite we want to arrive at a contradiction, let
$\alpha_1 \neq 0$ and $\alpha_2 = 0$ then $\alpha_1 x + 0 y= 0$ i.e.
$x = 0$ but we said that $x \neq 0$ so it must be that $\alpha_1 \neq 0$
and $\alpha_2 \neq 0$. Then we can write that
\begin{align*}
    x = -\frac{\alpha_2}{\alpha_1} y
\end{align*}
Also, we know that $\inprd{x,y} = 0$ since $x \perp y$ then
\begin{align*}
    \inprd{x,y} = \inprd{-\frac{\alpha_2}{\alpha_1} y, y}
    = -\frac{\alpha_2}{\alpha_1}\inprd{y,y} = 0
\end{align*}
But this implies that $y = 0$ and we said that $y \neq 0$, a contradiction.

Therefore must be that $\{x,y\}$ is a linearly independent set.

\item [(b)] Now, let $\{x_1, ..., x_m\}$ be a set of mutually orthogonal
non-zero vectors, we want to show they are a linearly independent set.

As before, suppose the opposite, we want to arrive at a contradiction.

If $\alpha_i \neq 0$ for only the $i$th vector then $x_i = 0$ and that
cannot happen.

So let us assume that $\alpha_i \neq 0$ for $1 \leq i \leq m$, if not all
$\alpha_i$s were different from zero the following still applies, then 
\begin{align*}
    \alpha_1 x_1 + ... + \alpha_n x_n &= 0\\
    x_1 &= -\frac{\alpha_2}{\alpha_1} x_2 - ... - \frac{\alpha_n}{\alpha_1} x_n
\end{align*}
Also we know that $\inprd{x_1, x_2} = 0$ then
\begin{align*}
    \inprd{x_1, x_2}
    &= \inprd{-\frac{\alpha_2}{\alpha_1} x_2 - ...
    - \frac{\alpha_n}{\alpha_1} x_n, x_2}\\
    &= -\frac{\alpha_2}{\alpha_1} \inprd{x_2, x_2} - ...
    - \frac{\alpha_n}{\alpha_1}\inprd{x_n, x_2}\\
    &= -\frac{\alpha_2}{\alpha_1} \inprd{x_2, x_2}\\
    &= 0
\end{align*}
Where we used that $\inprd{x_i, x_2} = 0$ for $1 \leq i \leq m$ and
$i \neq 2$.
But the last equality implies that $x_2 = 0$ which is a contradiction.

Therefore must be that the orthogonal nonzero vectors $\{x_1, ..., x_m\}$
are a linearly independent set.
\end{itemize}
\end{proof}

\cleardoublepage
\begin{proof}{\textbf{7}}
Let $\inprd{x,u} = \inprd{x,v}$ for all $x$, then
must also be true for $x = u$ i.e. $\inprd{u, u} = \inprd{u,v}$
and for $x = v$ i.e. $\inprd{v, u} = \inprd{v,v}$.\\
Now, let us compute $\inprd{u- v, u - v}$ as follows
\begin{align*}
    \inprd{u- v, u - v}
    &= \inprd{u, u} - \inprd{u,v} - \inprd{v, u} + \inprd{v, v}\\
    &= \inprd{u, v} - \inprd{u,v} - \inprd{v, u} + \inprd{v, u}\\
    &= 0
\end{align*}
But then this implies that $u - v = 0$. Therefore $u = v$.
\end{proof}
\begin{proof}{\textbf{8}}
Let us compute (9) knowing that $\|x\| = \sqrt{\inprd{x,x}}$ then 
\begin{align*}
    \frac{1}{4}(\|x + y\|^2 - \|x -y\|^2)
    &= \frac{1}{4}(\inprd{x + y, x + y} - \inprd{x - y, x - y})\\
    &= \frac{1}{4}(\inprd{x,x} + \inprd{x,y} + \inprd{y,x} + \inprd{y,y} -\\
    &\quad- (\inprd{x,x} + \inprd{x,-y} + \inprd{-y,x} + \inprd{-y,-y}))\\
    &= \frac{1}{4}(\inprd{x,x} + \inprd{x,y} + \inprd{y,x} + \inprd{y,y} -\\
    &\quad- \inprd{x,x} + \inprd{x,y} + \inprd{y,x} - \inprd{y,y})\\
    &= \frac{1}{4}(2\inprd{x,y} + 2\inprd{x,y})\\
    &= \inprd{x,y}
\end{align*}
Where we used that $\inprd{x,y} = \inprd{y,x}$ since we are in a real inner
product space.
\end{proof}

\cleardoublepage
\begin{proof}{\textbf{9}}
Let us compute (10) knowing that we are in a complex inner product space, then
\begin{align*}
    \frac{1}{4}(\|x + y\|^2 - \|x -y\|^2)
    &= \frac{1}{4}(\inprd{x + y, x + y} - \inprd{x - y, x - y})\\
    &= \frac{1}{4}(2\inprd{x,y} + 2\inprd{y,x})\\
    &= \frac{1}{4}(2\inprd{x,y} + 2\overline{\inprd{x,y}})\\
    &= \frac{1}{2}(\re{\inprd{x,y}} + i\im{\inprd{x,y}}
    + \re{\inprd{x,y}} - i\im{\inprd{x,y}})\\
    &= \re{\inprd{x,y}}
\end{align*}
Now, we compute the second equation as follows
\begin{align*}
    \frac{1}{4}(\|x + iy\|^2 - \|x -iy\|^2)
    &= \frac{1}{4}(\inprd{x + iy, x + iy} - \inprd{x - iy, x - iy})\\
    &= \frac{1}{4}(\inprd{x,x} + \inprd{x,iy} + \inprd{iy,x} + \inprd{iy,iy} -\\
    &\quad- (\inprd{x,x} + \inprd{x,-iy} + \inprd{-iy,x} + \inprd{-iy,-iy}))\\
    &= \frac{1}{4}(\inprd{x,x} - i\inprd{x,y} + i\inprd{y,x} + \inprd{y,y} -\\
    &\quad- \inprd{x,x} - i\inprd{x,y} + i\inprd{y,x} - \inprd{y,y})\\
    &= \frac{1}{2}(-i\inprd{x,y} + i\inprd{y,x})\\
    &= \frac{1}{2}(-i\inprd{x,y} + i\overline{\inprd{x,y}})\\
    &= \frac{1}{2}(-i\re{\inprd{x,y}} + \im{\inprd{x,y}} + i\re{\inprd{x,y}}
    + \im{\inprd{x,y}})\\
    &= \im{\inprd{x,y}}
\end{align*}
\end{proof}

\cleardoublepage
\begin{proof}{\textbf{10}}
Let $z_1$ and $z_2$ be complex numbers and let
$\inprd{z_1, z_2} = z_1\overline{z_2}$ we want to prove that in fact this
is a valid definition of an inner product.\\
So if this expression defines an inner product, it must have the four
properties of inner products.
Then for (IP1) let $z_3$ be also a complex number
\begin{align*}
    \inprd{z_1 + z_3, z_2} = (z_1 + z_3)\overline{z_2} 
    = z_1\overline{z_2} + z_3\overline{z_2}
    = \inprd{z_1, z_2} + \inprd{z_3, z_2}
\end{align*}
For (IP2) we see that
\begin{align*}
    \inprd{\alpha z_1, z_2} = \alpha z_1\overline{z_2}
    = \alpha\inprd{z_1, z_2}
\end{align*}
For (IP3) we see that
\begin{align*}
    \inprd{z_1, z_2} = z_1\overline{z_2}
    = \overline{\overline{z_1}z_2}
    = \overline{\inprd{z_2, z_1}}
\end{align*}
Finally for (IP4) we see that $\inprd{z_1, z_1} = z_1\overline{z_1} \geq 0$.
Also, let $\inprd{z_1, z_1} = 0$ then
\begin{align*}
    z_1\overline{z_1} =
    (x + iy)(x - iy) =
    x^2 + y^2 &= 0
\end{align*}
Then the only way for this equation to hold is that $x = y = 0$ and hence
$z_1 = 0$.
Now, if we suppose $z_1 = 0$ then
$\inprd{z_1, z_1} = z_1\overline{z_1} = 0$ and hence (IP4) is also
satisfied.\\
Therefore $\inprd{z_1, z_1} = z_1\overline{z_1}$ defines in fact an inner
product.\\
Also, to have orthogonality it must happen that
$\inprd{z_1, z_2} = z_1\overline{z_2} = 0$.
\end{proof}

\cleardoublepage
\begin{proof}{\textbf{11}}
Let $X$ be the vector space of all ordered pairs of complex numbers.
The norm defined as $\|x\| = |\xi_1| + |\xi_2|$ cannot be obtained from an
inner product. We prove this by showing that the norm does not satisfy the
parallelogram equality. Let us take $x = (-1, 1)$ and $y = (-1, -1)$
then
\begin{align*}
    2(\|x\|^2 + \|y\|^2) = 2((|-1| + |1|)^2 + (|-1| + |-1|)^2) = 16
\end{align*}
But 
\begin{align*}
    \|x+y\|^2 + \|x-y\|^2 = (|-1 -1| + |1 - 1|)^2 + (|-1 + 1| + |1 + 1|)^2
    = 8  
\end{align*}
Therefore the parallelogram equality is not satisfied and hence the norm
defined cannot be obtained from an inner product.
\end{proof}

\cleardoublepage
\begin{proof}{\textbf{13}}
The inner product defined in 3.1-5 for continuous functions states that
\begin{align*}
    \inprd{x,y} = \int_a^b x(t)\overline{y(t)}~dt
\end{align*}
Let us check that this definition satisfies (IP1) to (IP4).
We assume $x(t)$ and $y(t)$ complex-valued functions.\\
(IP1):
\begin{align*}
    \inprd{x + y,z} &= \int_a^b (x(t) + y(t))\overline{z(t)}~dt\\
    &= \int_a^b x(t)\overline{z(t)}~dt + \int_a^b y(t)\overline{z(t)}~dt\\
    &= \inprd{x, z} + \inprd{y, z}
\end{align*}
(IP2):
\begin{align*}
    \inprd{\alpha x,y} &= \int_a^b \alpha x(t)\overline{y(t)}~dt
    = \alpha\int_a^b x(t)\overline{y(t)}~dt
    = \alpha\inprd{x, y}
\end{align*}
(IP3):
\begin{align*}
    \inprd{x,y} &= \int_a^b x(t)\overline{y(t)}~dt
    = \int_a^b \overline{\overline{x(t)}y(t)}~dt
    = \overline{\inprd{y,x}}
\end{align*}
(IP4):
\begin{align*}
    \inprd{x, x} &= \int_a^b x(t)\overline{x(t)}~dt
    = \int_a^b |x(t)|^2~dt \geq 0
\end{align*}
Let now $\inprd{x, x} = 0$ then 
\begin{align*}
    \int_a^b |x(t)|^2~dt = 0
\end{align*}
Then must be that $x(t) = 0$ since $x(t)$ is continuous.
Conversely, if we suppose that $x(t) = 0$ we get that
\begin{align*}
    \inprd{x, x} &= \int_a^b 0\cdot 0~dt = 0
\end{align*}
Therefore the inner product defined in 3.1-5 satisfy the properties (IP1) to
(IP4).
\end{proof}

\cleardoublepage
\subsection*{3.2 - Further Properties of Inner Produet Spaces}
\begin{proof}{\textbf{2}}
An example of subspace of $l^2$ is the set sequences with finitely many entries
different from zero followed by zeroes i.e.
\begin{align*}
    Y = \{(x_1, x_2, ..., x_k, 0, 0, 0, ....): x_i \in \R, 0 \leq i \leq k\}
\end{align*}
Suppose $x = (x_1, x_2, ..., x_k, 0, 0, 0, ....) \in Y$ then 
$$\|x\|_2 = \bigg(\sum_{i=1}^\infty |x_i|^2\bigg)^{1/2}
= \bigg(\sum_{i=1}^k |x_i|^2\bigg)^{1/2} < \infty$$
So $Y \subset l^2$.\\
On the other hand, let $x' = (x_1', x_2', ..., x_k', 0, 0, 0, ....) \in Y$,
and $\alpha, \beta$ be scalars then 
\begin{align*}
    \alpha x + \beta x'
    &= \alpha (x_1, x_2, ..., x_k, 0, 0, 0, ....)
    + \beta (x_1', x_2', ..., x_k', 0, 0, 0, ....)\\
    &=(\alpha x_1, \alpha x_2, ..., \alpha x_k, 0, 0, 0, ....)
    + (\beta x_1', \beta x_2', ..., \beta x_k', 0, 0, 0, ....)\\
    &=(\alpha x_1 + \beta x_1', \alpha x_2 + \beta x_2', ...,
    \alpha x_k + \beta x_k', 0, 0, 0, ....)
\end{align*}
We see that $\alpha x_i + \beta x_i' \in \R$ since $\R$ is a vector space then
$Y$ is a subspace of $l^2$.
\end{proof}

\cleardoublepage
\begin{proof}{\textbf{3}}
Let $X$ be the inner product space consisting of the polynomial $x = 0$ and all
real polynomials in $t$ of degree not exceeding 2, considered for real
$t \in [a,b]$ with inner product defined as 
\begin{align*}
    \inprd{x,y} = \int_a^b x(t)y(t)~dt
\end{align*}
Let $(x_n) = (\alpha_n t^2 + \beta_n t + \gamma_n) \in X$ be a Cauchy sequence
on $X$ where $\alpha_n, \beta_n, \gamma_n \in \R$ then since this is a Cauchy
sequence, given $\epsilon > 0$ we can find $N \in \N$ such that when $n,m > N$
we have that
\begin{align*}
    \|x_n - x_m\| = \bigg(\int_a^b (x_n(t) - x_m(t))^2~dt\bigg)^{1/2}
    < \epsilon
\end{align*}
Hence
\begin{align*}
    \bigg(\int_a^b ((\alpha_n - \alpha_m)t^2 + (\beta_n - \beta_m)t + (\gamma_n - \gamma_m))^2~dt\bigg)^{1/2}
    < \epsilon
\end{align*}
For this integral to tend to 0 over the fixed interval $[a,b]$  since $t$ does
not depend on $n$ must happen that
$(\alpha_n - \alpha_m) \to 0$, $(\beta_n - \beta_m) \to 0$ and
$(\gamma_n - \gamma_m) \to 0$.
\\
This implies that $(\alpha_n)$, $(\beta_n)$ and $(\gamma_n)$ are Cauchy
sequences of real numbers which we know they converge,
let us say that they converge to $\alpha, \beta$ and $\gamma$ respectively.
\\
Now, let us define $x = \alpha t^2 + \beta t + \gamma$ then
\begin{align*}
    \|x_n - x\| &= \bigg(\int_a^b (x_n(t) - x(t))^2~dt\bigg)^{1/2}\\
    &= \bigg(\int_a^b ((\alpha_n - \alpha)t^2 + (\beta_n - \beta)t + (\gamma_n - \gamma))^2~dt\bigg)^{1/2}
\end{align*}
But we know that $\alpha_n \to \alpha$, $\beta_n \to \beta$ and
$\gamma_n \to \gamma$ therefore we have that
\begin{align*}
    \bigg(\int_a^b ((\alpha_n - \alpha)t^2 + (\beta_n - \beta)t + (\gamma_n - \gamma))^2~dt\bigg)^{1/2}
    \to 0 
\end{align*}
Or what is the same $x_n \to x$ and hence since $x_n$ was a Cauchy sequence
then $X$ is complete.
\\\\
Let now $Y$ be all $x \in X$ such that $x(a) = 0$. Let $y, y' \in Y$ and let
$\alpha, \beta$ be scalars. We know that when $t = a$ then $y(a) = y'(a) = 0$ then
for $z(t) = \alpha y(t) + \beta y'(t)$ we have that
$z(a) = \alpha y(a) + \beta y'(a) = 0 + 0 = 0$.
Therefore $Y$ is a subspace of $X$.

\cleardoublepage
Finally, let $Z$ be all $x\in X$ of degree 2. Let $z, z' \in Z$ and let
$\lambda, \rho$ be scalars. Then let us compute $\lambda z + \rho z'$ as
follows
\begin{align*}
    \lambda(\alpha t^2 + \beta t + \gamma) + \rho(\alpha' t^2 + \beta' t + \gamma')
    = (\lambda\alpha + \rho\alpha')t^2 + (\lambda\beta + \rho\beta')t + (\lambda\gamma + \rho\gamma')
\end{align*}
Where $\alpha, \alpha', \beta, \beta', \gamma$ and $\gamma'$ are not zero.
Then we see that if $\lambda = \rho = 0$ we get that $\lambda z + \rho z' = 0$
but the polynomial $z = 0$ is not included in $Z$ then $Z$ is not a subspace
of $X$.
\end{proof}
\begin{proof}{\textbf{4}}
Let $y \perp x_n$ and $x_n \to x$ then by the Continuity of the Inner Product
we have that
\begin{align*}
    \inprd{x_n, y} \to \inprd{x, y}
\end{align*}
This implies that
\begin{align*}
    \|\inprd{x_n, y} - \inprd{x, y}\| \to 0
\end{align*}
But we know that $\inprd{x_n, y} = 0$ because $y \perp x_n$ hence
\begin{align*}
    \|\inprd{x, y}\| \to 0
\end{align*}
And by the (N2) property for norms, we get that
\begin{align*}
    \inprd{x, y} = 0
\end{align*}
Therefore $x \perp y$.
\end{proof}

\cleardoublepage
\begin{proof}{\textbf{5}}
Let a sequence $(x_n)$ in an inner product space such that
$\|x_n\| \to \|x\|$ and $\inprd{x_n, x} \to \inprd{x, x}$.
Given that $\|x_n\| \to \|x\|$ then $\|x_n\|^2 \to \|x\|^2$ so given
$\epsilon^2/3 > 0$ there is $N \in \N$ such that when $n \geq N$ we see that
\begin{align*}
    |\|x_n\|^2 - \|x\|^2| = |\inprd{x_n, x_n} - \inprd{x,x}| < \epsilon^2/3
\end{align*}
On the other hand, since $\inprd{x_n, x} \to \inprd{x, x}$ given $\epsilon^2/3 > 0$
there is $N' \in \N$ such that when $n \geq N'$ we get that
\begin{align*}
    |\inprd{x_n, x} - \inprd{x,x}| &= |\inprd{x_n - x, x}| < \epsilon^2/3
\end{align*}
Also, let us note that
\begin{align*}
    \|x_n - x\|^2 &= |\inprd{x_n - x, x_n - x}|\\
    &= |\inprd{x_n - x, x_n} - \inprd{x_n - x, x}|\\
    &\leq |\inprd{x_n - x, x_n}| + |\inprd{x_n - x, x}|\\
    &\leq |\inprd{x_n, x_n} - \inprd{x, x_n}| + |\inprd{x_n - x, x}|\\
    &\leq |\inprd{x_n, x_n} - \inprd{x,x} + \inprd{x,x} - \inprd{x, x_n}| + |\inprd{x_n - x, x}|\\
    &\leq |\inprd{x_n, x_n} - \inprd{x,x}| + |\inprd{x, x - x_n}| + |\inprd{x_n - x, x}|
\end{align*}
Then if we take $M = \max(N, N')$ when $n \geq M$ we have that
\begin{align*}
    \|x_n - x\|^2 
    &\leq |\inprd{x_n, x_n} - \inprd{x,x}| + |\inprd{x, x - x_n}| + |\inprd{x_n - x, x}|\\
    & < \epsilon^2/3 + \epsilon^2/3 + \epsilon^2/3 = \epsilon^2 
\end{align*}
Therefore $\|x_n - x\| < \epsilon$ and hence $x_n \to x$.
\end{proof}

\cleardoublepage
\begin{proof}{\textbf{9}}
Let $V$ be the vector space of all continuous complex-valued functions on
$J = [a,b]$. Let $X_1 = (V, \|\cdot\|_\infty)$ where
$\|x\|_{\infty} = \max_{t \in J} |x(t)|$ and let $X_2 = (V, \|\cdot\|_2)$ where
$$\|x\|_2 = \sqrt{\inprd{x,x}} \qquad \inprd{x,y} = \int_{a}^b x(t)\overline{y(t)}~dt$$
Also, let us define $i: X_1 \to X_2$ be the identity mapping $x \to x$.\\
Suppose we take $x,y \in V$ then we see that
\begin{align*}
    \|x - y\|_2 &= \sqrt{\inprd{x - y,x - y}}\\
    &= \sqrt{\int_{a}^b (x(t) - y(t))(\overline{x(t) - y(t)})~dt}\\
    &= \sqrt{\int_{a}^b |x(t) - y(t)|^2~dt}
\end{align*}
But also we have that
\begin{align*}
    \sqrt{\int_{a}^b |x(t) - y(t)|^2~dt}
    &\leq \sqrt{\int_{a}^b \max_{t \in J}|x(t) - y(t)|^2~dt}
    &= \max_{t \in J}|x(t) - y(t)| \sqrt{b -a}
\end{align*}
So given $\epsilon > 0$ if we take $\delta = \epsilon/\sqrt{b - a}$ such that
\begin{align*}
    \|x - y\|_{\infty} = \max_{t \in J} |x(t) - y(t)| < \delta
    = \frac{\epsilon}{\sqrt{b - a}}
\end{align*}
We get that
\begin{align*}
    \|x - y\|_2 \leq \max_{t \in J}|x(t) - y(t)| \sqrt{b -a} < \epsilon
\end{align*}
Therefore the identity mapping $i$ is continuous.

\end{proof}

\cleardoublepage
\begin{proof}{\textbf{10}}
Let $T: X \to X$ be a bounded linear operator on a complex inner product space
$X$ and let $\inprd{Tx, x} = 0$ for all $x \in X$.
\\
Let $x = \alpha u + v \in X$ for some $u,v \in X$ and $\alpha \in \C$ then
\begin{align*}
    \inprd{Tx ,x} &= \inprd{\alpha Tu + Tv,\alpha u + v}\\
    &= \inprd{\alpha Tu, \alpha u + v} + \inprd{Tv,\alpha u + v}\\
    &= \alpha\overline{\alpha}\inprd{Tu, u} + \alpha\inprd{Tu, v}
    + \overline{\alpha}\inprd{Tv,u} + \inprd{Tv, v}\\
    &= \alpha\inprd{Tu, v} + \overline{\alpha}\inprd{Tv,u}
\end{align*}
Then taking $\alpha = 1$ we get that
\begin{align*}
    \inprd{Tu, v} + \inprd{Tv,u} &= 0\\
    \inprd{Tu, v} &= -\inprd{Tv,u} 
\end{align*}
But if we take $\alpha = i$ we have that $i\inprd{Tu, v} - i\inprd{Tv,u} = 0$
or
\begin{align*}
    \inprd{Tu, v} - \inprd{Tv,u} &= 0\\
    \inprd{Tu, v} &= \inprd{Tv,u}
\end{align*}
This implies that $\inprd{Tu, v} = 0$ for all $u,v \in X$ so in particular if
we select $v = Tu$ then $\inprd{Tu, Tu} = 0$ and hence by (IP4) we have that
$Tu = 0$. Therefore $T = 0$.
\\\\
Let us consider now the Euclidean plane, a real inner product space, and let
$T: \R^2 \to \R^2$ be defined by
\begin{align*}
    T\begin{bmatrix} x \\y \end{bmatrix}
    = \begin{bmatrix}
        \cos\pi/2 & -\sin\pi/2\\
        \sin\pi/2 & \cos\pi/2
    \end{bmatrix}\begin{bmatrix} x \\y \end{bmatrix}
    = \begin{bmatrix}
        0 & -1\\ 1 & 0
    \end{bmatrix}\begin{bmatrix} x \\y \end{bmatrix}
    = \begin{bmatrix} -y \\ x \end{bmatrix}
\end{align*}
Then let us compute the following inner product
\begin{align*}
    \inprd{T\begin{bmatrix} x \\y \end{bmatrix},
    \begin{bmatrix} x \\y \end{bmatrix}}
    &= \inprd{\begin{bmatrix} -y \\x \end{bmatrix},
    \begin{bmatrix} x \\y \end{bmatrix}}
    = -yx + xy = 0
\end{align*}
Therefore, we see that the inner product is 0 but $T \neq 0$ so above result
does not hold for real inner product spaces.
\end{proof}

\cleardoublepage
\subsection*{3.3 - Orthogonal Complements and Direct Sums}
\begin{proof}{\textbf{1}}
Let $H$ be a Hilbert space, $M \subset H$ a convex subset, and $(x_n)$ a
sequence in $M$ such that $\|x_n\| \to d$ where $d = \inf_{x \in M}\|x\|$.
\\
Let us consider $\overline{M}$ the closure of $M$ then $\overline{M}$ is
complete because $\overline{M}\subset H$, also, $(x_n) \in \overline{M}$.
\\
Let $n,m \in \N$ where $n \geq m$.
We see that $\frac{1}{2}x_n + \frac{1}{2}x_m \in M$ because $M$ is convex but
also we have that
\begin{align*}
    \bigg\|\frac{x_n + x_m}{2}\bigg\| \geq d = \inf_{x \in M} \|x\|
\end{align*}
Hence
\begin{align*}
    \|x_n + x_m\| \geq 2d
\end{align*}
On the other hand, by the parallelogram equality we get that
\begin{align*}
    \|x_n - x_m\|^2 &= - \|x_n + x_m\|^2 + 2(\|x_n\|^2 + \|x_m\|^2)\\
    &\leq -(2d)^2 + 2(\|x_n\|^2 + \|x_m\|^2)
\end{align*}
But since $\|x_n\| \to d$ as $n, m \to \infty$ then $\|x_n - x_m\| \to 0$.
Therefore $(x_n)$ is Cauchy, but this also implies that $(x_n)$ converges
because $\overline{M}$ is a subset of a Hilbert space.
\end{proof}

\cleardoublepage
\begin{proof}{\textbf{5}}
Let $X = \R^2$ then
\begin{itemize}
\item[(a)] Let $M = \{x\}$ where $x = (\xi_1, \xi_2) \neq 0$ then $M^\perp$ by
definition is 
\begin{align*}
    M^\perp = \{x'\in X~|~x'\perp M\} = \{x'\in X~|~\inprd{x',x} = 0\}
\end{align*}
Let $x' = (-\xi_2, \xi_1) \in X$ then we see that
\begin{align*}
    \inprd{x',x} = -\xi_1 \xi_2 + \xi_2 \xi_1 = 0
\end{align*}
So $(-\xi_2, \xi_1) \in M^\perp$ but $(-\alpha\xi_2, \alpha\xi_1)$ for
$\alpha \in \R$ is also in $M^\perp$.
\\
Therefore $M^\perp = \{(-\alpha\xi_2,\alpha\xi_1) \in X ~|~ \alpha \in \R\}$.
\item[(b)] Let $M = \{x_1, x_2\}$ where $x_1$ and $x_2$ are linearly
independent then
\begin{align*}
    M^\perp &= \{x'\in X~|~x'\perp M\}\\
    &= \{x'\in X~|~\inprd{x',x_1} = 0 \text{ and } \inprd{x',x_2} = 0\}
\end{align*}
Suppose $x' \in X$ satisfy $\inprd{x',x_1} = 0$ and $\inprd{x',x_2} = 0$,
then we have that
\begin{align*}
    \inprd{x',x_1} = \eta_1\xi_{11} + \eta_2\xi_{12} = 0
\end{align*}
and 
\begin{align*}
    \inprd{x',x_2} = \eta_1\xi_{21} + \eta_2\xi_{22} = 0
\end{align*}
Where we assumed that $x_1 = (\xi_{11}, \xi_{12})$, $x_2 = (\xi_{21}, \xi_{22})$
and $x' = (\eta_{1}, \eta_{2})$.
But since $x_1$ and $x_2$ are linearly independent then the equations
\begin{align*}
    c_{1}\xi_{11} + c_2\xi_{21} = 0\\
    c_{1}\xi_{12} + c_2\xi_{22} = 0
\end{align*}
have only one solution, $c_1 = c_2 = 0$.
\\
Therefore must be that $x' = 0$ and hence
\begin{align*}
    M^\perp &= \{0\}
\end{align*}
\end{itemize}
\end{proof}

\cleardoublepage
\begin{proof}{\textbf{6}}\\
Let $Y = \{x ~|~ x=(\xi_j) \in l^2, \xi_{2n} = 0, n \in \N\}$ we want to show
that $Y$ is a closed subspace of $l^2$.
\\
Let $(x_j) \subset Y$ and suppose it converges to some
$x = (\eta_j) \not\in Y$ we want to arrive at a contradiction, then
given $\epsilon > 0$ there is some $N\in \N$ such that when $n \geq N$
we have that $\|x_n - x\| < \epsilon$, also, we see that
\begin{align*}
    \|x_n - x\| &= \sqrt{\sum_{j=1}^\infty |\xi_j - \eta_j|^2}
    = \sqrt{\sum_{j=1}^\infty |\xi_{2j + 1} - \eta_{2j + 1}|^2
    + \sum_{j=1}^\infty |\eta_{2j}|^2}
\end{align*}
Since the first term on the RHS is non-negative then
\begin{align*}
    \|x_n - x\| \geq \sqrt{\sum_{j=1}^\infty |\eta_{2j}|^2}
\end{align*}
So we have a lower bound on $\|x_n - x\|$ but $\|x_n - x\| \to 0$, a
contradiction.
\\
Therefore, must be that $\eta_{2j} = 0$ so $x \in Y$ and hence $Y$ is closed.
\\\\
Given $x \in Y$, to find $Y^\perp$ we need to find all elements $z \in l^2$
such that $\inprd{x,z} = 0$. We see that
\begin{align*}
    \inprd{x,z} = \sum_{j=1}^\infty \xi_j\bar\eta_j
    = \sum_{j=1}^\infty \xi_{2j + 1}\bar\eta_{2j + 1}
    + \sum_{j=1}^\infty \xi_{2j}\bar\eta_{2j}
    = \sum_{j=1}^\infty \xi_{2j + 1}\bar\eta_{2j + 1}
\end{align*}
So for $\inprd{x,z} = 0$ then must be that $\bar\eta_{2j + 1} = 0$ because it
must hold for every $x \in Y$. Therefore
\begin{align*}
    Y^\perp = \{x ~|~ x=(\eta_j) \in l^2, \eta_{2n + 1} = 0, n \in \N\}
\end{align*}
\\
Finally, let $Y = \text{span}\{e_1, ..., e_n\} \subset l^2$ where
$e_j = (\delta_{jk})$.
\\
We see that an element $x \in Y$ can be written as
\begin{align*}
    x = \alpha_1 e_1 + ... + \alpha_n e_n
    = \bigg(\alpha_1, \alpha_2, ..., \alpha_n, 0, 0, ...\bigg)
\end{align*}
Where $\alpha_1, ..., \alpha_n \in \C$.
\\
Let then $z \in l^2$, we see that
\begin{align*}
    \inprd{x,z}
    = \alpha_1 \bar\eta_1 + \alpha_2 \bar\eta_2 + ... + \alpha_n \bar\eta_n 
    + 0\cdot \bar\eta_{n+1} + 0\cdot \bar\eta_{n+2} + ...
\end{align*}
So for $\inprd{x,z} = 0$ must be that $\eta_j = 0$ for $1\leq j \leq n$ and
hence
\begin{align*}
    Y^\perp = \{x ~|~x = (\xi_j)\in l^2
    \text{ such that }\xi_j = 0\text{ for } 1 \leq j \leq n\}
\end{align*}
\end{proof}

\cleardoublepage
\begin{proof}{\textbf{7}}\\
Let $A$ and $A \subset  B$ be nonempty subsets of an inner product space $X$
\begin{itemize}
\item [(a)] Let $a \in A$ and $a' \in A^\perp$ we see that $\inprd{a,a'} = 0$
hence $a \perp A^\perp$, also by definition we know that
\begin{align*}
    A^{\perp\perp} = \{x \in X ~|~ x \perp A^{\perp}\}
\end{align*}
Then $a \in A^{\perp\perp}$ and since $a$ was arbitrary we have that
\begin{align*}
    A \subset A^{\perp\perp}
\end{align*}
\item [(b)] Let $b \in B$, $a \in A$ and $b' \in B^\perp$, by definition
$b' \perp B$ hence $\inprd{b, b'} = 0$ but since $a$ is also in $B$ then
must be that $\inprd{a, b'} = 0$ and since $a$ and $b$ were arbitrary then
this implies that $b' \perp A$ and therefore
\begin{align*}
    B^{\perp} \subset A^\perp
\end{align*}
\item [(c)] From part (a) we have that $A \subset A^{\perp\perp}$ and from part
(b) taking $B = A^{\perp\perp}$ we have that
\begin{align*}
    B^{\perp} = A^{\perp\perp\perp} \subset A^\perp
\end{align*}
Also, from part (a) if we take $A$ as $A^\perp$ following the same steps we get
that
\begin{align*}
    A^\perp \subset A^{\perp\perp\perp}
\end{align*}
Therefore joining both results we get that $A^{\perp\perp\perp} = A^\perp$.
\end{itemize}
\end{proof}

\cleardoublepage
\begin{proof}{\textbf{9}}\\
Let $Y$ be a subspace of a Hilbert space $H$
\begin{itemize}
\item [$(\Rightarrow)$] Suppose $Y$ is closed in $H$ then by Lemma 3.3-6 we get
that $Y = Y^{\perp\perp}$.
\item [$(\Leftarrow)$] Suppose $Y = Y^{\perp\perp}$. We prove first that
$Y^\perp$ is closed in $H$.
\\
Let $(x_n) \in Y^\perp$ be a convergent sequence, that converges to $x \in H$.
\\
Let also, $y \in Y$ then we know that $\inprd{x_n, y} = 0$, then by continuity
of the inner product we have that
\begin{align*}
    \inprd{x_n,y} \to \inprd{x,y}
\end{align*}
but $\inprd{x_n,y} = 0$ then $\inprd{x,y} = 0$ and hence $x \in Y^\perp$.
So $Y^\perp$ is closed.
\\
We see that $Y^\perp$ is a subspace of $H$ since if $y_1', y_2' \in Y^\perp$
and $x \in Y$ then
\begin{align*}
    \inprd{\alpha y_1' + \beta y_2', x} =
    \alpha\inprd{y_1', x} + \beta\inprd{y_2', x} = 0
\end{align*}
So following the same logic we used to prove $Y^\perp$ is closed
assuming $Y$ is a subspace of $H$ we get that $Y^{\perp\perp}$ is closed 
but $Y = Y^{\perp\perp}$.
\\
Therefore $Y$ is closed.
\end{itemize}
\end{proof}

\cleardoublepage
\begin{proof}{\textbf{10}}\\
Let $M \neq \emptyset$ be any subset of a Hilbert space $H$.
\\
Let us notice that $M^{\perp\perp}$ is a closed subspace of $H$ so
$M^{\perp\perp}$ is a Hilbert space and hence without loss of generality we 
may assume that $M$ is a subset of $H= M^{\perp\perp}$.
\\
Let us take $x, y \in H = M^{\perp\perp}$ then $\inprd{x,y} \neq 0$ because
otherwise either $x$ or $y$ are not in $M^{\perp\perp}$ so the only way where
$\inprd{x,y} = 0$ is that $x = y = 0$ and hence $M^\perp = \{0\}$.
\\
Then by theorem 3.3-7 we have that the span of $M$ is dense in $M^{\perp\perp}$
but then
$$M \subset \overline{\text{span}(M)} = M^{\perp\perp}$$
Therefore $M^{\perp\perp}$ is the smallest closed subspace containing $M$.
\end{proof}

\cleardoublepage
\subsection*{3.4 - Orthonormal Sets and Sequences}
\begin{proof}{\textbf{1.}}
Let $X$ be an inner product space of finite dimension $n$, then, $X$ contains 
a lineraly independent set of $n$ vectors (a basis), let us call them
$B' = \{b'_1, ..., b'_n\}$.
\\
We want to define $B = \{b_1, ..., b_n\}$ such that the elements are
orthonormal. Let us define $b_1 \in X$ such that
\begin{align*}
    b_1 = \frac{1}{\|b'_1\|}\cdot b'_1
\end{align*}
We see that we can write $b'_2$ as
\begin{align*}
    b'_2 = \inprd{b'_2, b_1}b_1 + v_2
\end{align*}
where $v_2$ is non-zero since the elements of $B'$ are linearly independent,
then
\begin{align*}
    v_2 = b'_2 - \inprd{b'_2, b_1}b_1
\end{align*}
We see that
\begin{align*}
    \inprd{b_1, v_2} &= \inprd{b_1, b'_2 - \inprd{b'_2, b_1}b_1}\\
    &= \inprd{b_1, b'_2} - \inprd{b_1, \inprd{b'_2, b_1}b_1}\\
    &= \inprd{b_1, b'_2} - \overline{\inprd{b'_2, b_1}}\inprd{b_1, b_1}\\
    &= \inprd{b_1, b'_2} - \inprd{b_1, b'_2}\frac{1}{\|b'_1\|^2}\inprd{b'_1, b'_1}\\
    &= \inprd{b_1, b'_2} - \inprd{b_1, b'_2}\frac{\|b'_1\|^2}{\|b'_1\|^2}\\
    &= \inprd{b_1, b'_2} - \inprd{b_1, b'_2}\\
    &= 0
\end{align*}
Hence, we see that they are orthogonal. Then we define $b_2$ as
\begin{align*}
    b_2 = \frac{1}{\|v_2\|}\cdot v_2
\end{align*}
Now, to define $b_3$ we note that there is $v_3$ such that
\begin{align*}
    v_3 = b'_3 -\inprd{b'_3, b_1}b_1 - \inprd{b'_3, b_2}b_2
\end{align*}
We see that $v_3$ is non-zero since the elements of $B'$ are linearly
independent. Also, $v_3 \perp b_1$ and $v_3 \perp b_2$, then we define $b_3$ 
as
\begin{align*}
    b_3 = \frac{1}{\|v_3\|}\cdot v_3
\end{align*}
We can carry on this process $n$ times until we find
\begin{align*}
    v_n = b'_n - \sum_{i =1}^{n-1} \inprd{b'_n, b_i}b_i
\end{align*}
From which we define $b_n$ as
\begin{align*}
    b_n = \frac{1}{\|v_n\|}\cdot v_n
\end{align*}
Therefore we have found a set $B$ that is orthogonal pairwise, where each
element has norm equal to 1, so $B$ is an orthonormal set.
\\
Finally, because of Lemma 3.4-2 we know that $B$ is linearly independent so $B$
is a basis for $X$ as well.
\end{proof}

\cleardoublepage
\begin{proof}{\textbf{2.}}
Equation $(12^*)$ states that
\begin{align*}
    \sum_{k=1}^n |\inprd{x, e_k}|^2 \leq \|x\|^2
\end{align*}
Where $x \in \R^r$ and $r \geq n$.
We know from equation (9) that
\begin{align*}
    \|y\|^2 = \sum_{k=1}^n |\inprd{x, e_k}|^2
\end{align*}
So $y$ is the projection of $x$ onto a set of orthonormal vectors in $\R^n$
hence $(12^*)$ is telling us geometrically that the length of a projection
vector $y$ is less than or equal to the length of $x$.
\end{proof}

\cleardoublepage
\begin{proof}{\textbf{3.}}
Let $x, y \in X$ where $X$ is an inner product space.
Let us take an orthonormal set $\{e_1, ..., e_n\}$ such that
\begin{align*}
    e_1 = \frac{y}{\|y\|}
\end{align*}
Then we see that
\begin{align*}
    \sum_{k=1}^n |\inprd{x, e_k}|^2
    &= |\inprd{x, e_1}|^2 + \sum_{k=2}^n |\inprd{x, e_k}|^2\\
    &= \bigg|\bigg\langle x, \frac{y}{\|y\|}\bigg\rangle\bigg|^2
    + \sum_{k=2}^n |\inprd{x, e_k}|^2\\
    &= \frac{|\inprd{x, y}|^2}{\|y\|^2}
    + \sum_{k=2}^n |\inprd{x, e_k}|^2
\end{align*}
We note that
\begin{align*}
    \frac{|\inprd{x, y}|^2}{\|y\|^2} \leq \frac{|\inprd{x, y}|^2}{\|y\|^2}
    + \sum_{k=2}^n |\inprd{x, e_k}|^2
\end{align*}
Now, applying $(12^*)$ to $x$ we get that
\begin{align*}
    \frac{|\inprd{x, y}|^2}{\|y\|^2}
    \leq \frac{|\inprd{x, y}|^2}{\|y\|^2}
    + \sum_{k=2}^n |\inprd{x, e_k}|^2 
     = \sum_{k=1}^n |\inprd{x, e_k}|^2 \leq \|x\|^2
\end{align*}
Hence
\begin{align*}
    |\inprd{x, y}|^2 \leq \|x\|^2\|y\|^2
\end{align*}
And taking the square root we get Schwarz inequality
\begin{align*}
    |\inprd{x, y}| \leq \|x\|\|y\|
\end{align*}
\end{proof}

\cleardoublepage
\begin{proof}{\textbf{4.}}
Let us consider the orthonormal sequence
\begin{align*}
    e_1 &= (0, 1, 0, 0, ...)\\
    e_2 &= (0, 0, 1, 0, ...)\\
    e_3 &= (0, 0, 0, 1, ...)\\
    &\vdots
\end{align*}
And let us consider the sequence
\begin{align*}
    x = (1, 0, 0, 0, ...) \in l^2
\end{align*}
Then we see that
\begin{align*}
    \sum_{k=1}^\infty |\inprd{x, e_k}|^2 = 0
\end{align*}
But
\begin{align*}
    \|x\|^2 = 1^2 + 0 + 0 + ... = 1
\end{align*}
Therefore the strict inequality holds in this case
\begin{align*}
    \sum_{k=1}^\infty |\inprd{x, e_k}|^2 = 0 < \|x\|^2 = 1
\end{align*}
\end{proof}

\cleardoublepage
\begin{proof}{\textbf{5.}}
Let $(e_k)$ be an orthonormal sequence in an inner product space $X$, let also
$x \in X$ and
\begin{align*}
    y = \sum_{k=1}^n \alpha_k e_k \qquad \alpha_k = \inprd{x, e_k}
\end{align*}
We want to show that $x - y$ is orthogonal to the subspace
$Y_n = \text{span}\{e_1,..., e_n\}$.
Let us consider some $y' \in Y_n$ defined as
\begin{align*}
    y' = \sum_{k=1}^n \inprd{y', e_k} e_k 
\end{align*}
Then
\begin{align*}
    \inprd{x -y, y'} &= \inprd{x, y'}  - \inprd{y, y'}\\
    &= \inprd{x, \sum_{k=1}^n \inprd{y', e_k} e_k}
    - \inprd{\sum_{k=1}^n \inprd{x, e_k} e_k, \sum_{j=1}^n \inprd{y', e_j} e_j}\\
    &= \sum_{k=1}^n \overline{\inprd{y', e_k}}\inprd{x, e_k}
    - \sum_{k=1}^n\inprd{x, e_k}\inprd{e_k, \sum_{j=1}^n \inprd{y', e_j} e_j}\\
    &= \sum_{k=1}^n \overline{\inprd{y', e_k}}\inprd{x, e_k}
    - \sum_{k=1}^n\inprd{x, e_k}\sum_{j=1}^n \overline{\inprd{y', e_j}}\inprd{e_k, e_j}
\end{align*}
Since $\inprd{e_k, e_j} = 1$ only if $k = j$ otherwise it's 0 then we can write
that
\begin{align*}
    \inprd{x -y, y'}
    &= \sum_{k=1}^n \overline{\inprd{y', e_k}}\inprd{x, e_k}
    - \sum_{k=1}^n\inprd{x, e_k} \overline{\inprd{y', e_k}}
    = 0
\end{align*}
Therefore since $y'$ was arbitrary we see that $x -y$ is orthogonal to $Y_n$.
\end{proof}

\cleardoublepage
\begin{proof}{\textbf{6.}}
Let $\{e_1, ..., e_n\}$ be an orthonormal set in an inner product space $X$,
where $n$ is fixed. Let $x \in X$ and let $y = \beta_1 e_1 +... + \beta_ne_n$.
We want to show that $\|x -y\|$ is minimum if and only if
$\beta_j = \inprd{x,e_j}$ where $j = 1,..., n$.
\begin{itemize}
\item [$(\Rightarrow)$] Suppose $\|x - y\|$ is minimum then we note that
\begin{align*}
    x -y = x - \sum_{j=1}^n \beta_j e_j
    &= \bigg(x - \sum_{i=1}^n \inprd{x, e_i} e_i \bigg)
    + \bigg(\sum_{k=1}^n \inprd{x, e_k} e_k - \sum_{j=1}^n \beta_j e_j\bigg)
\end{align*}
Then we see that
\begin{align*}
    &\inprd{x - \sum_{i=1}^n \inprd{x, e_i} e_i
    , \sum_{k=1}^n \inprd{x, e_k}e_k - \sum_{j=1}^n\beta_j e_j}\\
    &= \inprd{x, \sum_{k=1}^n \inprd{x, e_k}e_k}
    - \inprd{x, \sum_{j=1}^n\beta_j e_j}
    - \inprd{\sum_{i=1}^n \inprd{x, e_i} e_i, \sum_{k=1}^n \inprd{x, e_k}e_k}\\
    &\quad+ \inprd{\sum_{i=1}^n \inprd{x, e_i} e_i,\sum_{j=1}^n\beta_j e_j}\\
    &= \sum_{k=1}^n \overline{\inprd{x, e_k}}\inprd{x, e_k}
    - \sum_{j=1}^n\overline{\beta_j}\inprd{x, e_j}
    - \sum_{i=1}^n\inprd{x, e_i}\sum_{k=1}^n\overline{\inprd{x, e_k}}\inprd{e_i, e_k}\\
    &\quad+ \sum_{i=1}^n\inprd{x, e_i}\sum_{j=1}^n\overline{\beta_j}\inprd{ e_i, e_j}\\
    &= \sum_{k=1}^n \overline{\inprd{x, e_k}}\inprd{x, e_k}
    - \sum_{j=1}^n\overline{\beta_j}\inprd{x, e_j}
    - \sum_{i=1}^n\inprd{x, e_i}\overline{\inprd{x, e_i}}
    + \sum_{i=1}^n\inprd{x, e_i}\overline{\beta_i}\\
    &= 0
\end{align*}
So the vectors are orthogonal, then by the Pythagorean realation we have that
\begin{align*}
    \|x - \sum_{j=1}^n \beta_j e_j\|^2
    &= \|x - \sum_{i=1}^n \inprd{x, e_i} e_i\|^2
    + \|\sum_{k=1}^n \inprd{x, e_k} e_k - \sum_{j=1}^n \beta_j e_j\|^2
\end{align*}
We see the last term is
\begin{align*}
    \|\sum_{k=1}^n (\inprd{x, e_k} - \beta_k)e_k\|^2
    &= \inprd{\sum_{k=1}^n (\inprd{x, e_k} - \beta_k)e_k,
    \sum_{j=1}^n (\inprd{x, e_j} - \beta_j)e_j}\\
    &=\sum_{k=1}^n(\inprd{x, e_k} - \beta_k)
    \sum_{j=1}^n (\inprd{x, e_j} - \beta_j)\inprd{e_k, e_j}\\
    &=\sum_{k=1}^n|\inprd{x, e_k} - \beta_k|^2
\end{align*}
Therefore since the first term, $\|x - \sum_{i=1}^n \inprd{x, e_i} e_i\|^2$,
is constant, and we said that $\|x -y\|^2$ is mininum then must be that
\begin{align*}
    \sum_{k=1}^n|\inprd{x, e_k} - \beta_k|^2 = 0
\end{align*}
And that can only happen if $\beta_k = \inprd{x, e_k}$ for all
$k \in \{1,...,n\}$.
\item [$(\Leftarrow)$] Suppose now that $\beta_j = \inprd{x, e_j}$ for all
$j \in \{1,...,n\}$ then 
\begin{align*}
    \|x - \sum_{j=1}^n \beta_j e_j\|^2
    &= \|x - \sum_{i=1}^n \inprd{x, e_i} e_i\|^2
    + \sum_{j=1}^n|\inprd{x, e_j} - \beta_j|^2
    = \|x - \sum_{i=1}^n \inprd{x, e_i} e_i\|^2
\end{align*}
Since $\|x - \sum_{i=1}^n \inprd{x, e_i} e_i\|^2$ is a constant then
$\|x -y\|$ cannot take a smaller value, i.e. it's the minimum when 
$\beta_j = \inprd{x, e_j}$.
\end{itemize}
\end{proof}

\cleardoublepage
\begin{proof}{\textbf{7.}}
Let $(e_j)$ be any orthonormal sequence in an inner product space $X$ and let
$x,y \in X$.
\\
Let us consider the Cauchy-Schwartz inequality over the sequences $(\inprd{x,e_k})$ 
and $(\inprd{y,e_k})$, then
\begin{align*}
    \sum_{k=1}^\infty |\inprd{x, e_k}\inprd{y, e_k}|
    \leq \sqrt{\sum_{j=1}^\infty|\inprd{x, e_j}|^2}
    \sqrt{\sum_{m=1}^\infty|\inprd{y, e_m}|^2}
\end{align*}
We know that $\sum_{j=1}^\infty|\inprd{x, e_j}|^2 \leq \|x\|^2$ because of (12)
hence
\begin{align*}
    \sum_{k=1}^\infty |\inprd{x, e_k}\inprd{y, e_k}|
    \leq \sqrt{\|x\|^2}\sqrt{\|y\|^2} = \|x\|\|y\|
\end{align*}
\end{proof}

\cleardoublepage
\begin{proof}{\textbf{8.}}
Let $x \in X$ where $X$ is an inner product space and let $n_m$ be the number
of $\inprd{x,e_k}$ such that $|\inprd{x,e_k}| > 1/m$.
\\
From equation (12) we know that
\begin{align*}
    \sum_{k=1}^\infty |\inprd{x,e_k}|^2 \leq \|x\|^2
\end{align*}
So we see that
\begin{align*}
     \frac{n_m}{m^2} < \sum_{k=1}^\infty |\inprd{x,e_k}|^2 \leq \|x\|^2
\end{align*}
Therefore, must be that
\begin{align*}
     n_m < m^2\|x\|^2
\end{align*}

\end{proof}

\cleardoublepage
\begin{proof}{\textbf{9.}}
Let the sequence $(x_0, x_1, x_2, ...)$ where $x_j(t) = t^j$ on the interval
$[-1,1]$ where 
\begin{align*}
    \inprd{x,y} = \int_{-1}^1 x(t)y(t)~dt
\end{align*}
To orthonormalize the first three terms of the sequence we follow the 
Gram-Schmidt process.
Let us compute first $\|x_0\|$ as follows
\begin{align*}
    \|x_0\| &= \sqrt{\inprd{x_0, x_0}} = \sqrt{\int_{-1}^1 x_0(t)x_0(t)~dt}
    = \sqrt{\int_{-1}^1 t^0t^0~dt} = \sqrt{\int_{-1}^1~dt} = \sqrt{2}
\end{align*}
So the first element of our orthonormalized sequence is
\begin{align*}
    e_0 = \frac{1}{\|x_0\|} x_0 = \frac{1}{\sqrt{2}}\cdot t^0 = \frac{1}{\sqrt{2}}
\end{align*}
Now we compute $v_1$ as 
\begin{align*}
    v_1 &= x_1 - \inprd{x_1, e_0}e_0\\
    &= t - \frac{1}{\sqrt{2}}\cdot\int_{-1}^1 t\cdot \frac{1}{\sqrt{2}}~dt\\
    &= t - \frac{1}{2}\int_{-1}^1 t ~dt\\
    &= t - \frac{1}{2}\bigg[\frac{t^2}{2}\bigg]_{-1}^1\\
    &= t
\end{align*}
So the norm of $v_1$ is
\begin{align*}
    \|v_1\| = \sqrt{\int_{-1}^1 t^2~dt} = \sqrt{\frac{2}{3}}
\end{align*}
Then, the second element of the sequence is given by
\begin{align*}
    e_1 &= \frac{1}{\|v_1\|} v_1 = \sqrt{\frac{3}{2}}t
\end{align*}
Finally, we compute $v_2$ as
\begin{align*}
    v_2 &= x_2 - \inprd{x_2, e_1}e_1 - \inprd{x_2, e_0}e_0\\
    &= t^2 - \sqrt{\frac{3}{2}}t\cdot\int_{-1}^1 t^2\cdot \sqrt{\frac{3}{2}}t~dt
    - \frac{1}{\sqrt{2}}\cdot\int_{-1}^1 t^2\cdot \frac{1}{\sqrt{2}}~dt\\
    &= t^2 - \frac{3}{2}t\cdot\int_{-1}^1 t^3~dt
    - \frac{1}{2}\int_{-1}^1 t^2~dt\\
    &= t^2 - \frac{1}{3}
\end{align*}
Hence
\begin{align*}
    \|v_2\|
    &= \sqrt{\int_{-1}^1 \bigg(t^2 - \frac{1}{3}\bigg)^2~dt}\\
    &= \sqrt{\int_{-1}^1 \bigg(t^4 - \frac{2t^2}{3} + \frac{1}{9}\bigg)~dt}\\
    &= \sqrt{\frac{2}{5} - \frac{4}{9} + \frac{2}{9}}\\
    &= \sqrt{\frac{8}{45}}
\end{align*}
so the third element of the sequence is
\begin{align*}
    e_2 = \frac{1}{\|v_2\|}v_2
    = \sqrt{\frac{45}{8}}\bigg(t^2 - \frac{1}{3}\bigg)
\end{align*}
Just as a small sanity check we compute $\|e_2\|$ as follows
\begin{align*}
    \|e_2\| &= \sqrt{\int_{-1}^1 \frac{45}{8}\bigg(t^2 - \frac{1}{3}\bigg)^2~dt}\\
    &= \sqrt{\frac{45}{8}}\sqrt{\int_{-1}^1 \bigg(t^2 - \frac{1}{3}\bigg)^2~dt}\\
    &= \sqrt{\frac{45}{8}}\|v_2\|\\
    &= \sqrt{\frac{45}{8}}\sqrt{\frac{8}{45}}\\
    &= 1
\end{align*}
\end{proof}

\cleardoublepage
\subsection*{3.5 - Series Related to Orthonormal Sequences and Sets}
\begin{proof}{\textbf{1.}}
Suppose (6) converges to $x$ i.e.
\begin{align*}
    \sum_{k=1}^\infty \alpha_ke_k = x
\end{align*}
Then by Theorem 3.5-2 we know that (7)
\begin{align*}
    \sum_{k=1}^\infty |\alpha_k|^2
\end{align*}
converges.
\\
Let us take the squared norm of $x$ then we have that
\begin{align*}
    \|x\|^2 = \bigg\|\sum_{k=1}^\infty \alpha_ke_k\bigg\|^2
\end{align*}
But since every $\alpha_ke_k$ is orthonormal then by the Pythagorean relation
we have that
\begin{align*}
    \bigg\|\sum_{k=1}^\infty \alpha_ke_k\bigg\|^2
    = \sum_{k =1}^\infty \|\alpha_ke_k\|^2
\end{align*}
Therefore
\begin{align*}
    \|x\|^2 = \bigg\|\sum_{k=1}^\infty \alpha_ke_k\bigg\|^2
    = \sum_{k =1}^\infty \|\alpha_ke_k\|^2
    = \sum_{k =1}^\infty |\alpha_k|^2\|e_k\|^2
    = \sum_{k =1}^\infty |\alpha_k|^2
\end{align*}
\end{proof}

\cleardoublepage
\begin{proof}{\textbf{3.}}
Let us consider the orthonormal sequence
\begin{align*}
    e_1 &= (0, 1, 0, 0, ...)\\
    e_2 &= (0, 0, 1, 0, ...)\\
    e_3 &= (0, 0, 0, 1, ...)\\
    &\vdots
\end{align*}
And let us consider the sequence
\begin{align*}
    x = (1, 0, 0, 0, ...) \in l^2
\end{align*}
Then we see that
\begin{align*}
    \sum_{k=1}^\infty \inprd{x, e_k}e_k = (0, 0, 0, 0, ...)
    \neq (1, 0, 0, 0, ...) = x
\end{align*}
\end{proof}

\cleardoublepage
\begin{proof}{\textbf{4.}}\\
Let $(x_j)\subset X$ such that $\|x_1\| + \|x_2\| + ...$ converges.
\\
We want to prove that $(s_n)$ is Cauchy where $s_n = \sum_{i=1}^n x_i$.
Taking $n', m' \in \N$ such that $n' < m'$, we see that 
\begin{align*}
    \bigg\|\sum_{i=1}^{n'} x_i - \sum_{j=1}^{m'} x_j\bigg\|
    = \bigg\|\sum_{i=n'+1}^{m'} x_i\bigg\|
    \leq \sum_{i=n'+1}^{m'} \|x_i\|
\end{align*}
On the other hand, since $\sum_i \|x_i\|$ converges then, it's Cauchy, so,
given $\epsilon > 0$ there is $N\in \N$ such that when $n,m \geq N$ we have that
\begin{align*}
    \bigg|\sum_{i=1}^n\|x_i\| - \sum_{j=1}^m\|x_j\|\bigg| < \epsilon
\end{align*}
We may assume that $n < m$, then
\begin{align*}
    \bigg|\sum_{i=n+1}^m\|x_i\|\bigg| = \sum_{i=n+1}^m\|x_i\|< \epsilon
\end{align*}
Hence, taking $n' = n$ and $m' = m$ we see that
\begin{align*}
    \bigg\|\sum_{i=1}^{n} x_i - \sum_{j=1}^{m} x_j\bigg\|
    = \bigg\|\sum_{i=n+1}^{m} x_i\bigg\|
    \leq \sum_{i=n+1}^{m} \|x_i\| < \epsilon
\end{align*}
Therefore $(s_n)$ is Cauchy.
\end{proof}

\cleardoublepage
\begin{proof}{\textbf{5.}}\\
From problem \textbf{4.} we know that if $X$ is an inner product space and
$\sum_{i} \|x_i\|$ converges then the sequence $(s_n)$ with
$s_n = \sum_{i=1}^n x_i$ is Cauchy.
\\
If instead we take $(s_n)$ in a Hilbert space $H$, the space is complete so
all Cauchy sequences converge, hence the sequence of partial sums $(s_n)$
converge to some limit $s$, so we have that
\begin{align*}
    \sum_{i=1}^\infty x_i = \lim_{n\to \infty}\sum_{i=1}^n x_i
    = \lim_{n\to \infty} s_n = s
\end{align*}
Which implies that $\sum_{i=1}^\infty x_i$ converges.
\end{proof}

\cleardoublepage
\begin{proof}{\textbf{6.}}\\
Let $(e_j)$ be an orthonormal sequence in Hilbert space $H$ and let 
\begin{align*}
    x = \sum_{j=1}^\infty \alpha_j e_j
    \qquad
    y = \sum_{j=1}^\infty \beta_j e_j
\end{align*}
Let us define the partial sums sequences
\begin{align*}
    x_n = \sum_{j=1}^n \alpha_j e_j
    \qquad
    y_n = \sum_{j=1}^n \beta_j e_j
\end{align*}
we see that $x_n \to x$ and $y_n \to y$ as $n \to \infty$.
\\
Then
\begin{align*}
    \inprd{x_n, y_n} 
    = \inprd{\sum_{j=1}^n \alpha_j e_j, \sum_{i=1}^n \beta_i e_i}
    = \sum_{j=1}^n \alpha_j\sum_{i=1}^n\overline{\beta_i}\inprd{e_j, e_i}
    = \sum_{j=1}^n \alpha_j\overline{\beta_j}
\end{align*}
Because $x_n \to x$, $y_n \to y$ and the continuity of the inner product we
have that $\inprd{x_n, y_n} \to \inprd{x, y}$. Therefore
\begin{align*}
    \inprd{x, y} = \lim_{n\to\infty}\inprd{x_n, y_n}
    = \lim_{n\to\infty} \sum_{j=1}^n \alpha_j\overline{\beta_j}
    = \sum_{j=1}^\infty \alpha_j\overline{\beta_j}
\end{align*}
\end{proof}

\cleardoublepage
\begin{proof}{\textbf{7.}}\\
Let $(e_k)$ be an orthonormal sequence in a Hilbert space $H$ and let $x\in H$.
We want to show that
$$y = \sum_{k=1}^\infty \inprd{x,e_k}e_k$$
exists in $H$ and that $x -y$ is orthogonal to every $e_k$.
\\
Let us consider the partial sum sequence
\begin{align*}
    y_n = \sum_{k=1}^n \inprd{x,e_k}e_k
\end{align*}
If we let $m < n$ then we see that
\begin{align*}
    \|y_n - y_m\|^2
    &= \bigg\|\sum_{k=1}^n \inprd{x,e_k}e_k - \sum_{j=1}^m \inprd{x,e_j}e_j\bigg\|^2\\
    &= \bigg\|\sum_{k=m+1}^n \inprd{x,e_k}e_k\bigg\|^2\\
    &= \sum_{k=m+1}^n \|\inprd{x,e_k}e_k\|^2\\
    &= \sum_{k=m+1}^n |\inprd{x,e_k}|^2
\end{align*}
On the other hand, by Bessel inequality we know that
$\sum_{k=1}^\infty |\inprd{x,e_k}|^2$ converges so the tail of the series
$\sum_{k=m+1}^n |\inprd{x,e_k}|^2$ must tend to 0 as $n,m \to \infty$.
\\
Therefore $\|y_n - y_m\|^2 \to 0$ as $n,m \to \infty$ and hence $(y_n)$ is
Cauchy and since $(y_n) \in H$ then $(y_n)$ converges. So 
\begin{align*}
    y = \lim_{n\to\infty}\sum_{k=1}^n \inprd{x,e_k}e_k = \sum_{k=1}^\infty \inprd{x,e_k}e_k
\end{align*}
exists in $H$.
\\
Let us consider now the following
\begin{align*}
    \inprd{x - y_n, e_k}
    &= \inprd{x - \sum_{j=1}^n \inprd{x,e_j}e_j, e_k}\\
    &= \inprd{x, e_k} - \inprd{\sum_{j=1}^n \inprd{x,e_j}e_j, e_k}\\
    &= \inprd{x, e_k} - \sum_{j=1}^n \inprd{x,e_j}\inprd{e_j, e_k}
\end{align*}
Now taking the limit as $n \to \infty$ we get that
\begin{align*}
    \lim_{n\to \infty}\inprd{x - y_n, e_k}
    &= \inprd{x, e_k} - \lim_{n\to \infty}\sum_{j=1}^n \inprd{x,e_j}\inprd{e_j, e_k}\\
    \inprd{x - y, e_k}
    &= \inprd{x, e_k} - \sum_{j=1}^\infty \inprd{x,e_j}\inprd{e_j, e_k}\\
    \inprd{x - y, e_k} &= \inprd{x, e_k} - \inprd{x,e_k}\\
    \inprd{x - y, e_k} &= 0
\end{align*}
Where we used the continuity of the inner product, to show that as $n\to\infty$
we have that $\inprd{x - y_n, e_k} \to \inprd{x - y, e_k}$.
\\
Therefore $x -y$ is orthogonal to every $e_k$.
\begin{align*}
    \inprd{\sum_{j=1}^\infty \inprd{x,e_j}e_j, e_k}
    &= \sum_{j=1}^\infty \inprd{x,e_j}\inprd{e_j, e_k}
\end{align*}
\end{proof}

\cleardoublepage
\begin{proof}{\textbf{8.}}\\
Let $(e_k)$ be an orthonormal sequence in a Hilbert space $H$ and let $M$ be
$M = \text{span}(e_k)$.
\begin{itemize}
\item [($\Rightarrow$)]
Let $x \in \bar{M}$ then $x$ can be written as 
\begin{align*}
    x = \sum_{k=1}^\infty \alpha_k e_k
\end{align*}
Where $\alpha_k$ are scalars. Since the summation converges to $x$ then by
Theorem 3.5-2 (b) we know the coefficients $\alpha_k$ are the Fourier 
coefficients $\inprd{x, e_k}$ so $x$ can be written as
\begin{align*}
    x = \sum_{k=1}^\infty \inprd{x, e_k}e_k
\end{align*}
\item [($\Leftarrow$)]
Let $x\in H$ such that it can be written as
\begin{align*}
    x = \sum_{k=1}^\infty \inprd{x, e_k}e_k
    = \sum_{k=1}^\infty \alpha_ke_k
\end{align*}
Then $x$ is an infinite linear combination of the elements of the orthonormal
sequence $(e_k)$ so must be that $x \in \bar{M}$.
\end{itemize}
\end{proof}

\cleardoublepage
\begin{proof}{\textbf{9.}}\\
Let $(e_n)$ and $(\tilde{e}_n)$ be orthonormal sequences in a Hilbert space $H$
and let $M_1 = \text{span}(e_n)$ and $M_2 = \text{span}(\tilde{e}_n)$.
\begin{itemize}
\item [($\Rightarrow$)] Let $\bar{M}_1 = \bar{M}_2$. By problem 8 we know
that if $x \in \bar{M}_1$ and $x \in H$ then it can be written as
\begin{align*}
    x = \sum_{m=1}^\infty \alpha_m e_{m} = \sum_{m=1}^\infty \inprd{x, e_m} e_{m}
\end{align*}
We note that $e_n \in \bar{M}_2$ then we can write it as
\begin{align*}
    e_n = \sum_{m=1}^\infty \inprd{e_n, \tilde{e}_{m}} \tilde{e}_{m}
    = \sum_{m=1}^\infty \alpha_{nm} \tilde{e}_{m}
\end{align*}
Where we defined $\alpha_{nm} = \inprd{e_n, \tilde{e}_{m}}$.\\
In the same way, since $\tilde{e}_n \in M_1$ then we can write it as
\begin{align*}
    \tilde{e}_n = \sum_{m=1}^\infty \inprd{\tilde{e}_n, e_{m}} e_{m}
    = \sum_{m=1}^\infty \overline{\inprd{e_{m}, \tilde{e}_n}} e_{m}
    = \sum_{m=1}^\infty \bar{\alpha}_{mn} e_{m}
\end{align*}
\item [($\Leftarrow$)] Let us suppose that
\begin{align*}
    e_n = \sum_{m=1}^\infty \alpha_{nm} \tilde{e}_{m}
    \quad\text{and}\quad
    \tilde{e}_n = \sum_{m=1}^\infty \bar{\alpha}_{mn} e_{m}
\end{align*}
Where $\alpha_{nm} = \inprd{e_n, \tilde{e}_{m}}$.\\
We see that $e_n$ and $\tilde{e}_n$ can be written in the form of equation (6)
then the orhtonormal sequence $(e_n)$ is both in $\bar{M}_1$ and $\bar{M}_2$
and the orhtonormal sequence $(\tilde{e}_n)$ is also in $\bar{M}_1$ and
$\bar{M}_2$.\\
Let us take $x \in \bar{M}_1$ then by problem 8 we know that $x$ can be written
as
\begin{align*}
    x &= \sum_{m=1}^\infty \inprd{x, e_m} e_{m}
\end{align*}
But since the orthonormal sequence $(e_n)$ is also in $\bar{M}_2$ then this
implies that $x \in \bar{M}_2$ and since $x$ is arbitrary then 
\begin{align*}
   \bar{M}_1 \subset \bar{M}_2
\end{align*}
The opposite can be proven if $x \in \bar{M}_2$ so we get that
\begin{align*}
   \bar{M}_1 = \bar{M}_2
\end{align*}
\end{itemize}
\end{proof}

\cleardoublepage
\subsection*{3.6 - Total Orthonormal Sets and Sequences}
\begin{proof}{\textbf{1.}}\\
Let $F$ be an orthonormal basis in an inner product space $X$. We know that
\begin{align*}
    \overline{\Span{F}} = X
\end{align*}
so if $x \in X$ can be represented as a linear combination of elements of $F$
then $x \in \Span{F}$ but it could happen that $x \in X$ but it cannot be
represented as a finite linear combination of elements in $F$.
\\
For example, let us consider $X = l^2$ then if we take $F$ as
\begin{align*}
    F = \{
        (1,0, 0, ...), 
        (0, 1, 0, ...),
        (0, 0, 1, ...),
        ...
    \}
\end{align*}
We see that $\overline{\Span{F}} = l^2$, also, the sequence 
$(1, 1/2, 1/3, ...) \in l^2$ but is not in $\Span{F}$ since it's infinite.
\end{proof}

\begin{proof}{\textbf{2.}}\\
Suppose the orthogonal dimension of a hilbert space $H$ is finite, then there
is a total orthonormal set $T$ with finite number of elements in it,
let us suppose this number is $n$, but by Lemma 3.4-2 we know that orthonormal
sets are linealy independent so inside $H$ we have a set of $n$ linearly
independent elements.
\\
Let us suppose that there is, for example, one more element that is linearly
independent with respect to the $n$ previous elements, we want to arrive at 
a contradiction. So, in total we have $n + 1$ linearly independent elements,
then, we can build an orthonormal set $F$ of $n + 1$ elements using the
Gram-Schmidt process, since this new orthonormal set was built from the
original total orthonormal set $T$ then also, must be that $\Span{F}$ is
dense in $H$, but for this to happen by Lemma 3.3-7 must be that
$\Span{F}^\perp = \{0\}$ which implies that $\Span{T}^\perp \neq \{0\}$ in the
first place.
\\
This is a contradiction and must be that there cannot be more than $n$
linearly independent sets if the orthogonal dimension of $H$ is $n$.
\\
Therefore, the vector space $H$ is of dimension $n$.
\\
Now, suppose $H$ is considered as a vector space is of dimension $n$ (finite),
then there is a linearly independent set with $n$ elements in it.
We can obtain an orthonormal set of $n$ elements by applying the Gram-Schmidt
process and we observe that all elements of $H$ are in the span of the
orthonormal set, so the orthonormal set is total, hence, the orthogonal
dimension of $H$ is $n$.
\end{proof}

\cleardoublepage
\begin{proof}{\textbf{4.}}\\
Let us consider first the case of a complex inner product space.
From Parseval relation we note that
\begin{align*} \|x - y\|^2 &= \sum_{k} |\inprd{x - y, e_k}|^2 \end{align*}
Hence, from the LHS we see that
\begin{align*} 
    \|x - y\|^2 &= \inprd{x - y, x - y}
    = \inprd{x, x} - \inprd{x,y} - \inprd{y, x} + \inprd{y,y}
\end{align*}
And, from the RHS we have that
\begin{align*}
    \sum_{k} |\inprd{x - y, e_k}|^2
    &= \sum_{k} \inprd{x - y, e_k}\overline{\inprd{x - y, e_k}}\\
    &= \sum_{k} (\inprd{x, e_k}-\inprd{y, e_k})(\overline{\inprd{x, e_k}} - \overline{\inprd{y,e_k}})\\
    &= \sum_{k} \inprd{x, e_k}\overline{\inprd{x, e_k}} - \inprd{y, e_k}\overline{\inprd{x, e_k}} - \inprd{x, e_k}\overline{\inprd{y,e_k}} + \inprd{y, e_k}\overline{\inprd{y,e_k}}\\
    &= \sum_{k} |\inprd{x, e_k}|^2 - \inprd{y, e_k}\overline{\inprd{x, e_k}} - \inprd{x, e_k}\overline{\inprd{y,e_k}} + |\inprd{y, e_k}|^2\\
    &= \inprd{x,x} + \inprd{y,y} -\sum_{k} \inprd{y, e_k}\overline{\inprd{x, e_k}} + \inprd{x, e_k}\overline{\inprd{y,e_k}}
\end{align*}
Joining both results we see that
\begin{align*}
    \inprd{x,y} + \inprd{y, x}
    &= \sum_{k} \inprd{y, e_k}\overline{\inprd{x, e_k}} + \inprd{x, e_k}\overline{\inprd{y,e_k}}
\end{align*}
Now, let us consider
\begin{align*}
    \|x - iy\|^2 &= \sum_{k} |\inprd{x - iy, e_k}|^2
\end{align*}
From the LHS we see that
\begin{align*}
    \|x - iy\|^2 &= \inprd{x - iy, x - iy}\\
    &= \inprd{x , x} - \inprd{x,iy} - \inprd{iy, x} + \inprd{iy, iy}\\
    &= \inprd{x , x} + i\inprd{x, y} - i\inprd{y, x} + \inprd{y, y}
\end{align*}
And, from the RHS we have that
\begin{align*}
    \sum_{k} |\inprd{x - iy, e_k}|^2
    &= \sum_{k} \inprd{x - iy, e_k}\overline{\inprd{x - iy, e_k}}\\
    &= \sum_{k} (\inprd{x, e_k}-i\inprd{y, e_k})(\overline{\inprd{x, e_k}}
    + i\overline{\inprd{y,e_k}})\\
    &= \sum_{k} \inprd{x, e_k}\overline{\inprd{x, e_k}}
    - i\inprd{y, e_k}\overline{\inprd{x, e_k}}
    + i\inprd{x, e_k}\overline{\inprd{y,e_k}} 
    + \inprd{y, e_k}\overline{\inprd{y,e_k}}\\
    &= \sum_{k} |\inprd{x, e_k}|^2
    - i\inprd{y, e_k}\overline{\inprd{x, e_k}}
    + i\inprd{x, e_k}\overline{\inprd{y,e_k}} 
    + |\inprd{y, e_k}|^2\\
    &= \inprd{x,x} + \inprd{y,y} + i\sum_{k}
    \inprd{x, e_k}\overline{\inprd{y,e_k}}
    -\inprd{y, e_k}\overline{\inprd{x, e_k}}
\end{align*}
Again, joining both results we see that
\begin{align*}
    \inprd{x, y} - \inprd{y, x} &= \sum_{k}
    \inprd{x, e_k}\overline{\inprd{y,e_k}}
    -\inprd{y, e_k}\overline{\inprd{x, e_k}}
\end{align*}
Adding both result equations we get that
\begin{align*}
    &\inprd{x,y} + \inprd{y, x}
    + \inprd{x, y} - \inprd{y, x}\\
    &\quad= \sum_{k}
    \inprd{y, e_k}\overline{\inprd{x, e_k}}
    + \inprd{x, e_k}\overline{\inprd{y,e_k}}
    + \sum_{k}
    \inprd{x, e_k}\overline{\inprd{y,e_k}}
    -\inprd{y, e_k}\overline{\inprd{x, e_k}}
\end{align*}
Hence
\begin{align*}
    \inprd{x,y} = \sum_{k} \inprd{x, e_k}\overline{\inprd{y,e_k}}
\end{align*}

Let us consider now the case of a real inner product space.
From the polarization identity we know that
\begin{align*}
    \inprd{x,y} = \frac{1}{2}(\|x\|^2 + \|y\|^2 - \|x-y\|^2)
\end{align*}
Hence
\begin{align*}
    \inprd{x,y} &= \frac{1}{2}(
    \inprd{x,x} + \inprd{y,y} - \inprd{x,x} - \inprd{y,y}
    + \sum_{k}
    \inprd{y, e_k}\inprd{x, e_k}
    + \inprd{x, e_k}\inprd{y,e_k})\\
    &= \frac{1}{2}(
    2\sum_{k}\inprd{x, e_k}\inprd{y,e_k})\\
    &= \sum_{k} \inprd{x, e_k}\inprd{y,e_k}
\end{align*}
Where we used the previous result we got for $\|x-y\|^2$ and that
$\inprd{x, e_k} = \inprd{e_k, x}$ in a real inner product space.
\end{proof}

\cleardoublepage
\begin{proof}{\textbf{5.}}\\
Let $(e_\kappa)$, with $\kappa \in I$ be an orthonormal family in a Hilbert
space $H$.
\begin{itemize}
\item [$(\Rightarrow)$] Let us suppose $(e_\kappa)$ is total.
Since Parseval relation
$$\sum_k |\inprd{x,e_k}|^2 = \|x\|^2$$
Holds for any total orthonormal set $M$ in a Hilbert space $H$ and we derived 
the equation 
\begin{align*}
    \inprd{x,y} = \sum_k\inprd{x,e_k}\overline{\inprd{y,e_k}}
\end{align*}
Using the Parseval relation then it must also hold for the orthonormal family
$(e_\kappa)$.
\item [$(\Leftarrow)$] Let us suppose that for every $x,y \in H$ the following 
equation holds
\begin{align*}
    \inprd{x,y} = \sum_\kappa\inprd{x,e_\kappa}\overline{\inprd{y,e_\kappa}}
\end{align*}
For the orthonormal family $(e_\kappa)$. Taking $y = x$ we get that 
\begin{align*}
    \|x\|^2 = \inprd{x,x}
    = \sum_\kappa\inprd{x,e_\kappa}\overline{\inprd{x,e_\kappa}}
    = \sum_\kappa |\inprd{x,e_\kappa}|^2
\end{align*}
Then Parseval relation also holds and by Theorem 3.6-3 then must be that
the orthonormal family $(e_\kappa)$ is total.
\end{itemize}
\end{proof}

\cleardoublepage
\begin{proof}{\textbf{6.}}\\
Let $H$ be a separable Hilbert space and $M = \{x_1, x_2, x_3, ...\}$ a
countable dense subset of $H$.
Then, let us select
\begin{align*}
    e_1 = \frac{x_1}{\|x_1\|}
\end{align*}
To select $e_2$ we check if $x_2$ is in $\Span{e_1}$, if it is there then we
drop it, if it's not then we define $v_2 = x_2 - \inprd{x_2, e_1} e_1$
and we select
\begin{align*}
    e_2 = \frac{v_2}{\|v_2\|}
\end{align*}
To select $e_3$ we check if $x_3$ is in $\Span{\{e_1, e_2\}}$, if it is there
then we drop it, if it's not then we define
$v_3 = x_3 - \inprd{x_3, e_1} e_1 - \inprd{x_3, e_2} e_2$
and we select
\begin{align*}
    e_3 = \frac{v_3}{\|v_3\|}
\end{align*}
i.e. we follow the Gram-Schmidt process dropping some elements if necessary, to
get an orthonormal sequence $(e_j) \subseteq M$.
\\
Since the elements we dropped are in some span then they can be written as a
linear combination of elements of $(e_j)$, so $M \subseteq \Span{e_j}$
but also we know that $M$ is a dense subset of $H$ so must be that $\Span{e_j}$
is also dense in $H$.
\\
Therefore $(e_j)$ is a total orthonormal sequence obtained from $M$.
\end{proof}

\cleardoublepage
\begin{proof}{\textbf{7.}}\\
Let $H$ be a separable Hilbert space, then there is $M$ a countable dense
subset in $H$. As we showed before, using the Gram-Schmidt process we can find
a total orthonormal sequence in $H$.
\end{proof}

\cleardoublepage
\begin{proof}{\textbf{8.}}\\
Let $F = (e_j)$ be an orthonormal sequence in a separable Hilbert space $H$.\\
We know there exists $D = \{d_1, d_2, d_3, ...\}$, a countable dense subset
of $H$, because $H$ is separable.\\
Let us take $d_1$, if $d_1$ is linearly independent from the rest of the
elements in $D$ we keep it.\\
Next, we compute
\begin{align*}
    v_1 = d_1 - \sum_{k=1}^{\infty}\inprd{d_1, e_k}e_k
\end{align*}
and
\begin{align*}
    x_1 = \frac{v_1}{\|v_1\|}
\end{align*}
i.e. we follow the Gram-Schmidt process to orthonormalize $d_1$.\\
Finally, we add $x_1$ to a new orthogonal sequence $\tilde{F}$ defined as
\begin{align*}
    \tilde{F} = \{x_1, x_2, ...,  e_1, e_2, ...\}
\end{align*}
To get each $x_i$ we follow the same process for each $d_i$.
If $d_i$ is not linearly independent we drop it and continue.\\
Following this process we see that $\tilde{F}$ is an orthonormal sequence, 
and since we started from a countable dense subset $D$, then $\Span{\tilde{F}}$
must also be dense in $H$, since we only dropped the elements that were
linearly dependent to the other elements.\\
Therefore $\tilde{F}$ is a total orthonormal sequence that contains $F$.
\end{proof}

\cleardoublepage
\begin{proof}{\textbf{9.}}\\
Let $M$ be a total set in an inner product space $X$ and
let $\inprd{v,x} = \inprd{w,x}$ for all $x \in M$. We note that
\begin{align*}
    \inprd{v,x} &= \inprd{w,x}\\
    \inprd{v,x} - \inprd{w,x} &= 0\\
    \inprd{v - w, x} &= 0
\end{align*}
This implies that $v - w\perp x$ but because Theorem 3.6-2, there is no element
in $X$ which is orthogonal to every element in $M$ so must be that $v -w = 0$
and therefore $v = w$.
\end{proof}

\cleardoublepage
\begin{proof}{\textbf{10.}}\\
Let $M$ be a subset of a Hilbert space $H$ and let $v, w \in H$. Suppose that
$\inprd{v,x} = \inprd{w,x}$ for all $x \in M$ implies that $v = w$.
\\
If $M$ is not total then there is some $y \in H$, $y \neq 0$ such that
$y \perp M$, we want to arrive at a contradiction.
\\
Then $\inprd{y, x} = 0$, but also we know that 
$\inprd{0, x} = 0$ so $\inprd{y, x} = \inprd{0, x}$ hence $y = 0$ but $y$ was
different from 0, a contradiction.
\\ 
Therefore $M$ is total in $H$.
\end{proof}

\cleardoublepage
\subsection*{3.7 - Legendre, Hermite and Laguerre Polynomials}
\begin{proof}{\textbf{1.}}\\
The Legendre differential equation is
\begin{align*}
    (1 - t^2)P_n'' - 2tP_n' + n(n + 1)P_n = 0
\end{align*}
We note that applying the product rule to $[(1 - t^2)P_n']'$ gives us
\begin{align*}
    [(1 - t^2)P_n']' = (1 - t^2)'P_n' + (1 - t^2)P_n''
    = (1 - t^2)P_n'' - 2tP_n'
\end{align*}
Then we can write the Legendre differential equation as
\begin{align*}
    [(1 - t^2)P_n']' = -n(n + 1)P_n 
\end{align*}
Multiplying this equation by $P_m$, where $m \neq n$, gives us
\begin{align*}
    P_m[(1 - t^2)P_n']' = -n(n + 1)P_nP_m
\end{align*}
The corresponding equation for $P_m$ multiplied by $-P_n$ is
\begin{align*}
    -P_n[(1 - t^2)P_m']' = m(m + 1)P_mP_n
\end{align*}
Then, adding the equation and integrating from -1 to 1 we see that
\begin{align*}
    \int_{-1}^1 P_m[(1 - t^2)P_n']' -P_n[(1 - t^2)P_m']'~dt
    &= \int_{-1}^1 m(m + 1)P_mP_n -n(n + 1)P_nP_m ~dt\\
    &= [m(m + 1) -n(n + 1)]\int_{-1}^1 P_nP_m~dt
\end{align*}
Let us integrate $\int_{-1}^1 P_m[(1 - t^2)P_n']'~dt$ by parts 
\begin{align*}
    \int_{-1}^1 P_m[(1 - t^2)P_n']'~dt
    &= P_m[(1 - t^2)P_n'] \bigg|_{-1}^1 - \int_{-1}^1 P_m'[(1 - t^2)P_n']~dt\\
    &= 0 - \int_{-1}^1 (1 - t^2)P_n'P_m'~dt
\end{align*}
In the same way for $\int_{-1}^1 P_n[(1 - t^2)P_m']'~dt$ we get that
\begin{align*}
    \int_{-1}^1 P_n[(1 - t^2)P_m']'~dt
    &= P_n[(1 - t^2)P_m'] \bigg|_{-1}^1 - \int_{-1}^1 P_n'[(1 - t^2)P_m']~dt\\
    &= 0 - \int_{-1}^1 (1 - t^2)P_n'P_m'~dt
\end{align*}
Therefore
\begin{align*}
    &\int_{-1}^1 P_m[(1 - t^2)P_n']' -P_n[(1 - t^2)P_m']'~dt\\
    &\quad= - \int_{-1}^1 (1 - t^2)P_n'P_m'~dt + \int_{-1}^1 (1 - t^2)P_n'P_m'~dt\\
    &\quad= 0
\end{align*}
So
\begin{align*}
    \inprd{P_n, P_m} = \int_{-1}^1 P_nP_m~dt = 0
\end{align*}
Hence $(P_n)$ is an orthogonal sequence in the space $L^2[-1,1]$.

\end{proof}

\cleardoublepage
\begin{proof}{\textbf{2.}}\\
The binomial theorem states that
\begin{align*}
    (a + b)^n = \sum_{j=0}^n \frac{n!}{j! (n -j)!} a^{n-j} b^j
\end{align*}
So in this case we can apply it to $(t^2 - 1)^n$ as follows
\begin{align*}
    (t^2 - 1)^n = \sum_{j=0}^n \frac{n!}{j! (n -j)!} t^{2n-2j} (-1)^j
\end{align*}
Let us take the $n$th derivative of $t^{2n - 2j}$, we get that
\begin{align*}
    \dv[n]{t}t^{2n-2j}
    &= (2n - 2j) \dv[n-1]{t}t^{2n - 2j - 1}\\ 
    &= (2n - 2j)\cdot(2n - 2j - 1) \dv[n-2]{t}t^{2n - 2j - 2}\\
    &\quad\vdots\\
    &= [(2n - 2j)\cdot(2n - 2j - 1)\cdot ... \cdot(2n - 2j - n)]t^{2n - 2j - n}\\
    &= \frac{(2n - 2j)!}{(n - 2j)!}t^{n - 2j}
\end{align*}
Hence
\begin{align*}
    \dv[n]{t}(t^2 - 1)^n = \sum_{j=0}^n (-1)^j \frac{n!}{j! (n -j)!}
    \frac{(2n - 2j)!}{(n - 2j)!}t^{n - 2j}
\end{align*}
Therefore
\begin{align*}
    P_n(t) &= \frac{1}{2^n n!}\dv[n]{t}[(t^2 - 1)^n] \\
    &= \frac{1}{2^n n!}\sum_{j=0}^n (-1)^j \frac{n!}{j! (n -j)!}
    \frac{(2n - 2j)!}{(n - 2j)!}t^{n - 2j}\\
    &= \sum_{j=0}^n (-1)^j \frac{(2n - 2j)!}{2^n j! (n -j)! (n - 2j)!} t^{n - 2j}
\end{align*}



\end{proof}

\cleardoublepage
\begin{proof}{\textbf{3.}}\\
Let us consider the LHS as a binomial series
\begin{align*}
    \frac{1}{\sqrt{1 - 2tw + w^2}}
    = (1 - 2tw + w^2)^{-1/2}
    = \sum_{k = 0}^\infty \binom{-1/2}{k} (- 2tw + w^2)^k
\end{align*}
Now applying the binomial theorem we get that
\begin{align*}
    \frac{1}{\sqrt{1 - 2tw + w^2}}
    &= \sum_{k = 0}^\infty \binom{-1/2}{k} \sum_{j = 0}^k \binom{k}{j}(- 2tw)^{k-j} w^{2j}\\
    &= \sum_{k = 0}^\infty \binom{-1/2}{k}
    \sum_{j = 0}^k (-1)^{k-j}\binom{k}{j}2^{k-j} t^{k-j} w^{2j + k - j}\\
    &= \sum_{k = 0}^\infty \sum_{j = 0}^k
    (-1)^{k-j}\frac{(-1/2) (-1/2 - 1) ... (-1/2 - k + 1)}{k!}
    \frac{k!}{j!(k - j)!}2^{k-j} t^{k-j} w^{j + k}\\
    &= \sum_{k = 0}^\infty \sum_{j = 0}^k
    (-1)^{k-j}\frac{2^k(-1/2) (-1/2 - 1) ... (-1/2 - k + 1)}{2^kj!(k - j)!}
    2^{k-j} t^{k-j} w^{j + k}\\
    &= \sum_{k = 0}^\infty \sum_{j = 0}^k
    (-1)^{k-j}\frac{2^k\prod_{i=0}^{k-1} -(2i + 1)/2}{2^j j!(k - j)!}
    t^{k-j} w^{j + k}\\
    &= \sum_{k = 0}^\infty \sum_{j = 0}^k
    (-1)^{k-j}\frac{\prod_{i=0}^{k-1} -(2i + 1)}{2^j j!(k - j)!}
    t^{k-j} w^{j + k}\\
    &= \sum_{k = 0}^\infty \sum_{j = 0}^k
    (-1)^{k-j}\frac{(-1)^{k}\prod_{i=0}^{k-1}(2i + 1)}{2^j j!(k - j)!}
    t^{k-j} w^{j + k}\\
    &= \sum_{k = 0}^\infty \sum_{j = 0}^k
    (-1)^{2k-j}\frac{2k !}{2^{k+ j} k! j!(k - j)!}
    t^{k-j} w^{j + k}\\
    &= \sum_{k = 0}^\infty \sum_{j = 0}^k
    (-1)^{j}\frac{2k !}{2^{k+ j} k! j!(k - j)!}
    t^{k-j} w^{j + k}
\end{align*}
Where we used in the last step that $(-1)^{2k -j} = (-1)^{2k}(-1)^{-j}
= 1 \cdot (-1)^{j}$.\\
Applying the formula $\sum_{k=0}^\infty \sum_{j=0}^k C_{j,k}t^{k+j}
= \sum_{k=0}^\infty \sum_{j=0}^{k/2} C_{j, k-j}t^{k}$ we get that
\begin{align*}
    \frac{1}{\sqrt{1 - 2tw + w^2}}
    &= \sum_{k = 0}^\infty \sum_{j = 0}^{k/2}
    (-1)^{j}\frac{(2(k -j))!}{2^{k - j+ j} (k-j)! j!(k - j - j)!}
    t^{(k - j) - j} w^{j + k - j}\\
    &= \sum_{k = 0}^\infty \sum_{j = 0}^{k/2}
    (-1)^{j}\frac{(2k - 2j)!}{2^{k} (k-j)! j!(k - 2j)!}
    t^{k -2j} w^{k}\\
    &= \sum_{k = 0}^\infty P_k(t) w^{k}
\end{align*}
\end{proof}

\cleardoublepage
\begin{proof}{\textbf{4.}}\\
Let us consider 
\begin{align*}
    \frac{1}{r} = \frac{1}{\sqrt{r_1^2 + r_2^2 - 2r_1r_2\cos\theta}}
\end{align*}
but let us write it as
\begin{align*}
    \frac{1}{\sqrt{r_1^2 + r_2^2 - 2r_1r_2\cos\theta}}
    = \frac{1}{r_2}\frac{1}{\sqrt{1 + (r_1/r_2)^2 - 2(r_1/r_2)\cos\theta}}
\end{align*}
Then taking $w = r_1/r_2$ and $t = \cos\theta$ the RHS has the form 
$1 /r_2 \sqrt{1 - 2tw + w^2}$ and hence using the result of the previous
problem we get that
\begin{align*}
    &\frac{1}{r_2}\frac{1}{\sqrt{1 - 2(r_1/r_2)\cos\theta + (r_1/r_2)^2}} =\\
    &\quad= \frac{1}{r_2}\frac{1}{\sqrt{1 - 2wt + w^2}}\\
    &\quad= \frac{1}{r_2}\sum_{n=0}^\infty P_n(t)w^n\\
    &\quad= \frac{1}{r_2}\sum_{n=0}^\infty P_n(\cos\theta)
    \bigg(\frac{r_1}{r_2}\bigg)^n
\end{align*}
\end{proof}

\cleardoublepage
\begin{proof}{\textbf{5.}}\\
Let us consider $x(t) = \sum_{k=0}^\infty c_k t^k$, we note that
\begin{align*}
    x(t)' &= \sum_{k=1}^\infty c_k kt^{k-1} \\
    x(t)'' &= \sum_{k=2}^\infty c_k k(k-1)t^{k-2}
\end{align*}
Replacing this into Legendre's differential equation gives us
\begin{align*}
    (1 - t^2)x(t)'' - 2tx(t)' + n(n+1)x(t) &= 0\\
    %
    (1 - t^2)\sum_{k=2}^\infty c_k k(k-1) t^{k-2}
    - 2t\sum_{k=1}^\infty c_k kt^{k-1}
    + n(n+1)\sum_{k=0}^\infty c_k t^{k} &= 0\\
    %
    \sum_{k=2}^\infty c_k k(k-1) t^{k-2}
    - \sum_{k=2}^\infty c_k k(k-1) t^{k}
    - \sum_{k=1}^\infty 2c_kk t^{k}
    + \sum_{k=0}^\infty n(n+1)c_k t^{k} &= 0\\
    %
    \sum_{k=0}^\infty c_{k + 2} (k + 1)(k + 2) t^{k}
    - \sum_{k=2}^\infty c_k k(k-1) t^{k}
    - \sum_{k=1}^\infty 2c_kk t^{k}
    + \sum_{k=0}^\infty n(n+1)c_k t^{k} &= 0
\end{align*}
In the second term for $k = 0$ and $1$ the terms are 0, so we can set the sum to
start at $k = 0$, the same applies for the third term, hence
\begin{align*}
    \sum_{k=0}^\infty [c_{k + 2} (k + 1)(k + 2) - c_k k(k-1) 
    - 2c_kk + n(n+1)c_k]t^{k} &= 0
\end{align*}
Since the sum is equal to 0 then must be that each term is 0, so we get that
\begin{align*}
    c_{k + 2} (k + 1)(k + 2) &= [k(k-1) 
    + 2k - n(n+1)]c_k\\
    c_{k + 2} &= \frac{k(k-1) 
    + 2k - n(n+1)}{(k + 1)(k + 2)}c_k
\end{align*}
We see that
\begin{align*}
    c_2 &= -\frac{n(n+1)}{2}c_0 & k &= 0\\
    c_3 &= \frac{2 - n(n + 1)}{6}c_1 = -\frac{(n-1)(n+2)}{6}c_1 & k &= 1\\
    c_4 &= \frac{6 -n(n+1)}{12}c_2 = -\frac{(n -2)(n + 3)}{12}c_2 = & k &= 2\\
        &= \frac{(n -2)(n + 3)}{12}\frac{n(n+1)}{2}c_0 = 
        \frac{(n -2)(n + 3)n(n+1)}{24}c_0 &   &  \\
    &... 
\end{align*}
Then we can write $x(t)$ as
\begin{align*}
    x(t) &= c_0\bigg[
        1 - \frac{n(n+1)}{2!}t^2 + \frac{(n -2)(n + 3)n(n+1)}{4!}t^4 + ...
    \bigg] + \\
    &+ c_1\bigg[
        t - \frac{(n-1)(n+2)}{3!}t^3
        + \frac{(n - 3)(n - 1)(n + 2)(n + 4)}{5!}t^5 + ...
    \bigg]
\end{align*}
or defining $x_1(t)$ an $x_2(t)$ as
\begin{align*}
    x_1(t) &= 1 - \frac{n(n+1)}{2!}t^2 + \frac{(n -2)(n + 3)n(n+1)}{4!}t^4 + ...\\
    x_2(t) &= t - \frac{(n-1)(n+2)}{3!}t^3
    + \frac{(n - 3)(n - 1)(n + 2)(n + 4)}{5!}t^5 + ...
\end{align*}
We can write it as $x(t) = c_0 x_1(t) + c_1 x_2(t)$.\\
Finally, we note that when $n$ is even, then the term containing $t^{n+2}$ a
$x_1(t)$ becomes 0, and since the coefficients are recursive, all subsequent
coefficients become 0. The same thing happens when $n$ is odd, so $x_2(t)$
becomes also a polynomial.\\
On the other hand, if we assume that 
\begin{align*}
    c_{n-2j} = (-1)^j \frac{(2n - 2j)!}{2^n j! (n-j)! (n -2j)!}
\end{align*}
Then $c_{n- 2(j-1)} = c_{n -2j +2}$ must be
\begin{align*}
    c_{n-2j+ 2} = (-1)^{j-1} \frac{(2n - 2j + 2)!}{2^n (j-1)! (n-j+1)! (n -2j + 2)!}
\end{align*}
But also by the recurrence relation, setting $k=n - 2j$ we have that
\begin{align*}
    c_{n - 2j + 2} &= \frac{(n-2j)(n-2j-1) + 2(n - 2j) - n(n+1)}
    {(n - 2j + 1)(n - 2j + 2)}c_{n-2j}\\
    &= (-1)^j\frac{(n-2j)(n-2j-1) + 2(n - 2j) - n(n+1)}
    {(n - 2j + 1)(n - 2j + 2)}\frac{(2n - 2j)!}{2^n j! (n-j)! (n -2j)!}\\
    &= (-1)^j[(n-2j)(n-2j + 1) - n(n+1)]
    \frac{(2n - 2j)!}{2^n j! (n-j)! (n -2j + 2)!}\\
    &= (-1)^j[n^2 -2jn + n -2jn + 4j^2 - 2j - n(n+1)]
    \frac{(2n - 2j)!}{2^n j! (n-j)! (n -2j + 2)!}\\
    &= (-1)^j[-4jn + 4j^2 - 2j]
    \frac{(2n - 2j)!}{2^n j! (n-j)! (n -2j + 2)!}\\
    &= (-1)^j[-2j(2n - 2j + 1)]
    \frac{(2n - 2j)!}{2^n j! (n-j)! (n -2j + 2)!}\\
    &= (-1)^{j-1}\frac{2j!}{(j-1)!}
    \frac{(2n - 2j + 1)!}{2^n j! (n-j)! (n -2j + 2)!}\\
    &= (-1)^{j-1}\frac{2(n - j + 1)}{(n -j + 1)}
    \frac{(2n - 2j + 1)!}{2^n (j-1)! (n-j)! (n -2j + 2)!}\\
    &= (-1)^{j-1}\frac{(2n - 2j + 2)!}{2^n (j-1)! (n-j + 1)! (n -2j + 2)!}
\end{align*}
Therefore the coefficients of $x_1$ and $x_2$ match the coefficients of the
Legendre polynomials, if we set correctly the $c_0$ and $c_1$.
\end{proof}

\cleardoublepage
\begin{proof}{\textbf{6.}}\\
Let us consider the LHS as
\begin{align*}
    e^{2wt - w^2} = \sum_{n=0}^\infty \frac{1}{n!}(2wt - w^2)^n
\end{align*}
Applying the binomial theorem we get that
\begin{align*}
    e^{2wt - w^2}
    &= \sum_{n=0}^\infty \frac{1}{n!} 
    \sum_{k=0}^n\binom{n}{k}(2wt)^{n-j} (-w^2)^{j}\\
    &= \sum_{n=0}^\infty \frac{1}{n!} 
    \sum_{j=0}^n\frac{n!}{j!(n-j)!}2^{n-j}w^{n-j}t^{n-j}(-1)^jw^{2j}\\
    &= \sum_{n=0}^\infty \frac{1}{n!}\bigg[
    n!\sum_{j=0}^n(-1)^j\frac{2^{n-j}}{j!(n-j)!}t^{n-j}w^{n + j}\bigg]
\end{align*}
Also, we know that $\sum_{n=0}^\infty \sum_{j=0}^n C_{j,n}t^{n+j}
= \sum_{n=0}^\infty \sum_{j=0}^{[n/2]} C_{j, n-j}t^{n}$.\\
Therefore
\begin{align*}
    e^{2wt - w^2}
    &= \sum_{n=0}^\infty \frac{1}{n!}\bigg[
    n!\sum_{j=0}^{[n/2]} (-1)^j\frac{2^{n-2j}}{j!(n-2j)!}t^{n-2j}w^{n}\bigg]
    = \sum_{n=0}^\infty \frac{1}{n!} H_n(t) w^{n}
\end{align*}
\end{proof}

\cleardoublepage
\begin{proof}{\textbf{7.}}\\
We know that
\begin{align*}
    H_n(t) = (-1)^{n}e^{t^2}\dv[n]{t}(e^{-t^2})
\end{align*}
So we see that
\begin{align*}
    H_n'(t) &= (-1)^n \bigg[
        2te^{t^2}\dv[n]{t}(e^{-t^2})
        + e^{t^2}\dv[n+1]{t}(e^{-t^2})
    \bigg]\\
    -H_n'(t) &= -2t(-1)^ne^{t^2}\dv[n]{t}(e^{-t^2})
    - (-1)^ne^{t^2}\dv[n+1]{t}(e^{-t^2})\\
    -H_n'(t) &= -2t(-1)^ne^{t^2}\dv[n]{t}(e^{-t^2})
    + (-1)^{n + 1}e^{t^2}\dv[n+1]{t}(e^{-t^2})
\end{align*}
Now, replacing the definition of $H_n(t)$ and that 
\begin{align*}
    H_{n+1}(t) &= (-1)^{n+1}e^{t^2}\dv[n+1]{t}(e^{-t^2})
\end{align*}
We get that
\begin{align*}
    -H_n'(t) &= -2tH_n(t) + H_{n+1}(t)
\end{align*}
Therefore
\begin{align*}
    H_{n+1}(t) &= 2tH_n(t) - H_{n}'(t)
\end{align*}
\end{proof}

\cleardoublepage
\begin{proof}{\textbf{8.}}\\
Let us differentiate the generating function as follows
\begin{align*}
    \dv{t}(e^{2wt - w^2}) = 2we^{2wt - w^2}
    = \sum_{n=1}^\infty \frac{1}{n!}H_n'(t) w^n
\end{align*}
The first term of the sum becomes 0 since $H_0(t) = 1$, so the sum starts at 1.
Then we have that
\begin{align*}
    2w\sum_{n=0}^\infty \frac{1}{n!}H_n(t) w^n
    &= \sum_{n=1}^\infty \frac{1}{n!}H_n'(t) w^n\\
    2w\sum_{n=1}^\infty \frac{1}{(n-1)!}H_{n-1}(t) w^{n-1}
    &= \sum_{n=1}^\infty \frac{1}{n!}H_n'(t) w^n\\
    2\sum_{n=1}^\infty \frac{1}{n!} n H_{n-1}(t) w^{n}
    &= \sum_{n=1}^\infty \frac{1}{n!}H_n'(t) w^n\\
    \sum_{n=1}^\infty \frac{1}{n!} [2n H_{n-1}(t) - H_n'(t)] w^n &= 0
\end{align*}
Where in the second step we changed the LHS sum to start from 1. Since the sum
is equal to 0 then must be that each term is 0, hence
\begin{align*}
    H_n'(t) = 2n H_{n-1}(t)
\end{align*}
The Hermite differential equation states that
\begin{align*}
    H_n'' - 2tH_n' + 2nH_n = 0
\end{align*}
From the equation we derived and the result of Problem 7 with $n-1$ 
we see that
\begin{align*}
    H_n'' = 2n H_{n-1}' = 2n(2tH_{n-1} - H_{n})
\end{align*}
Now, replacing in the Hermite differential equation we get that
\begin{align*}
    H_n'' - 2tH_n' + 2nH_n
    &= 2n(2tH_{n-1} - H_{n}) - 2t(2n H_{n-1}) + 2nH_n\\
    &= 4ntH_{n-1} - 2nH_{n} - 4tn H_{n-1} + 2nH_n\\
    &= 0
\end{align*}
Therefore $H_n$ satisfies the Hermite differential equation.
\end{proof}

\cleardoublepage
\begin{proof}{\textbf{9.}}\\
In the asymptotic case the differential equation $y'' + (2n +1 -t^2)y = 0$
becomes $y'' = t^2y$ for which we know that the solution is $y(t) = e^{-t^2/2}$
so we may assume that the solution to the original differential equation is of
the form $y(t) = h(t)e^{-t^2/2}$, replacing this solution into the differential
equation we have that
\begin{align*}
    \dv[2]{y}{t} + (2n +1 -t^2)y &= 0\\
    \dv[2]{t}(h(t)e^{-t^2/2}) + (2n +1 -t^2)h(t)e^{-t^2/2} &= 0\\
    e^{-t^2/2}(h(t)'' - 2th(t)' + (t^2 - 1)h(t))
    + (2n +1 -t^2)h(t)e^{-t^2/2} &= 0\\
    h(t)'' - 2th(t)' + (t^2 - 1)h(t)
    + 2nh(t) - (t^2 - 1)h(t) &= 0\\
    h(t)'' - 2th(t)' + 2nh(t) &= 0
\end{align*}
We see that the differential equation is of the form of the Hermite
differential equation, so $H_n$ is a solution to the equation.\\
Therefore $H_n(t)e^{-t^2/2}$ is a solution to the original differential
equation.
\end{proof}

\cleardoublepage
\begin{proof}{\textbf{10.}}\\
Let us take the derivative of $e^{-t^2}H_n'$ as follows
\begin{align*}
    (e^{-t^2}H_n')' = -2te^{-t^2}H_n' + e^{-t^2}H_n''
\end{align*}
Also, from Problem 8 we know that $H_n' = 2nH_{n-1}$ so
$$H_n'' = 2nH_{n-1}' = 2n(2tH_{n-1} - H_n)$$
Where we used that $H_{n-1}' = 2tH_{n-1} - H_n$ from Problem 7.\\
Hence replacing
\begin{align*}
    (e^{-t^2}H_n')'
    &= -4nte^{-t^2}H_{n-1} + 4nte^{-t^2}H_{n-1} - 2ne^{-t^2}H_n\\
    &= -2ne^{-t^2}H_n
\end{align*}
From the equation we just derived if we let $n \neq m$ and following the same
method as in Problem 1, we can write that
\begin{align*}
    H_m(e^{-t^2}H_n')' - H_n(e^{-t^2}H_m')'
    &= -2ne^{-t^2}H_nH_m + 2me^{-t^2}H_nH_m\\
    &= 2e^{-t^2}(m - n)H_nH_m
\end{align*}
The (7a) functions state that
\begin{align*}
    e_n = A_n e^{-t^2/2}H_n
\end{align*}
where $A_n = 1/(2^n n! \sqrt{\pi})^{1/2}$.\\
We want to check that the functions are orthogonal on $\R$ so we need
to integrate $e_n e_m$ from $-\infty$ to $\infty$ 
\begin{align*}
    A_nA_m\int_{-\infty}^{\infty}  e^{-t^2/2}H_ne^{-t^2/2}H_m~dt
    &= A_nA_m\int_{-\infty}^{\infty}  e^{-t^2}H_n H_m~dt
\end{align*}
We note that this is the same as integrating
$H_m(e^{-t^2}H_n')' - H_n(e^{-t^2}H_m')'$ times a multiplicative constant, so
integrating the first term by parts we see that
\begin{align*}
    \int_{-\infty}^{\infty}  (e^{-t^2}H_n')'H_m~dt
    &= e^{-t^2}H_n'H_m\bigg|_{t=-\infty}^{t=\infty}
    - \int_{-\infty}^{\infty}  e^{-t^2}H_n'H_m'~dt\\
    &= 0 - \int_{-\infty}^{\infty}  e^{-t^2}H_n'H_m'~dt
\end{align*}
In hte same way, integrating by parts the second term we get that
\begin{align*}
    -\int_{-\infty}^{\infty}  (e^{-t^2}H_m')'H_n~dt
    &= -e^{-t^2}H_m'H_n\bigg|_{t=-\infty}^{t=\infty}
    + \int_{-\infty}^{\infty}  e^{-t^2}H_m'H_n'~dt\\
    &= 0 + \int_{-\infty}^{\infty}  e^{-t^2}H_n'H_m'~dt
\end{align*}
Therefore
\begin{align*}
    \int_{-\infty}^\infty H_m(e^{-t^2}H_n')' - H_n(e^{-t^2}H_m')'~dt = 
    2(m-n)\int_{-\infty}^\infty e^{-t^2}H_nH_m~dt = 0
\end{align*}
Hence
\begin{align*}
    \inprd{e_n, e_m}
    &= A_nA_m\int_{-\infty}^{\infty}  e^{-t^2}H_n H_m~dt
    = 0
\end{align*}
Which impliest that the functions $e_n(t)$ are orthogonal.
\end{proof}

\cleardoublepage
\begin{proof}{\textbf{11.}}\\
Let us consider the LHS of the equation as follows
\begin{align*}
    \frac{1}{1 - w} e^{-tw/(1-w)}
    &= \frac{1}{1 - w}\sum_{n=0}^\infty \frac{1}{n!}\bigg(-\frac{tw}{1-w}\bigg)^n\\
    &= \sum_{n=0}^\infty \frac{1}{n!}\frac{(-1)^nw^n}{(1-w)^{n+1}}t^n
\end{align*}
Using the series expansion for $1/(1-w)^{n+1}$ we get that
\begin{align*}
    \frac{1}{1 - w} e^{-tw/(1-w)}
    &= \sum_{n=0}^\infty \frac{1}{n!}(-1)^nw^n t^n
    \sum_{k=0}^\infty\binom{k+n}{k} w^{k}\\
    &= \sum_{n=0}^\infty \sum_{k=0}^\infty
    \frac{(-1)^n}{n!}\binom{k+n}{k} w^{n+k}t^n\\
    &= \sum_{n=0}^\infty \sum_{j=n}^\infty
    \frac{(-1)^n}{n!}\binom{j}{n} w^{j}t^n
\end{align*}
Where in the last step we changed variables so $j = k + n$ and hence the sum
starts at $j = n$.\\
If $n = 0$ then $j$ can go from 0 to $\infty$, if $n = 1$ then $j$ can go from
$1$ to $\infty$ and so on. If we swap the summation signs for $j = 0$ then
$n$ can only be 0, for $j = 1$ then $n$ can only be 0 or 1 and so on, then
must be that
\begin{align*}
    \frac{1}{1 - w} e^{-tw/(1-w)}
    &= \sum_{j=0}^\infty \sum_{n=0}^j
    \frac{(-1)^n}{n!}\binom{j}{n} w^{j}t^n
\end{align*}
Or renaming variables $n \to j$ and $j \to n$ we get that
\begin{align*}
    \frac{1}{1 - w} e^{-tw/(1-w)}
    &= \sum_{n=0}^\infty \sum_{j=0}^n
    \frac{(-1)^j}{j!}\binom{n}{j} t^j w^{n}
\end{align*}
But from equation (10c) we know that
\begin{align*}
    L_n(t) = \sum_{j=0}^n \frac{(-1)^j}{j!} \binom{n}{j}t^j
\end{align*}
Therefore
\begin{align*}
    \frac{1}{1 - w} e^{-tw/(1-w)}
    &= \sum_{n=0}^\infty L_n(t) w^{n}
\end{align*}
\end{proof}

\cleardoublepage
\begin{proof}{\textbf{12.}}\\
\begin{itemize}
\item [(a)] Let us differentiate $\psi$ with respect to $w$
\begin{align*}
    \dv{\psi(t,w)}{w} &= \dv{w}(\frac{1}{1 - w}e^{-tw/(1-w)})\\
    &= \frac{1}{(1 - w)^2}e^{-tw/(1-w)}
    - \frac{1}{(1-w)}\frac{te^{-tw/(1-w)}}{(1 - w)^2}\\
    &= \frac{1}{1-w}\sum_{n=0}^\infty L_n w^n
    - \frac{t}{(1-w)^2}\sum_{n=0}^\infty L_{n} w^n
\end{align*}
From the RHS we have that
\begin{align*}
    \dv{w}(\sum_{n=0}^\infty L_n w^n) &= \sum_{n=1}^\infty n L_n w^{n-1}
    = \sum_{n=0}^\infty (n + 1) L_{n + 1} w^{n}
\end{align*}
So
\begin{align*}
    (1 - w)\sum_{n=0}^\infty L_n w^n
    - t\sum_{n=0}^\infty L_{n} w^n
    &= (1 - w)^2\sum_{n=0}^\infty (n + 1) L_{n + 1} w^{n}
\end{align*}
Hence
\begin{align*}
    &\sum_{n=0}^\infty L_n w^n - \sum_{n=0}^\infty L_n w^{n+1}
    - t\sum_{n=0}^\infty L_{n} w^n
    =\\
    &= \sum_{n=0}^\infty (n + 1) L_{n + 1} w^{n}
    -2\sum_{n=0}^\infty (n + 1) L_{n + 1} w^{n+1}
    +\sum_{n=0}^\infty (n + 1) L_{n + 1} w^{n+2}\\[6pt]
    %
    &(1 - t)\sum_{n=0}^\infty L_{n} w^n
    - \sum_{n=1}^\infty L_{n-1} w^{n}
    =\\
    &= \sum_{n=0}^\infty (n + 1) L_{n + 1} w^{n}
    -2\sum_{n=1}^\infty nL_{n} w^{n}
    +\sum_{n=2}^\infty (n - 1) L_{n-1} w^{n}
\end{align*}
If we make every summation to start from $n=2$ then 
\begin{align*}
    &A + (1 - t)\sum_{n=2}^\infty L_{n} w^n
    - \sum_{n=2}^\infty L_{n-1} w^{n}
    =\\
    &= B + \sum_{n=2}^\infty (n + 1) L_{n + 1} w^{n}
    -2\sum_{n=2}^\infty nL_{n} w^{n}
    +\sum_{n=2}^\infty (n - 1) L_{n-1} w^{n}
\end{align*}
Where
\begin{align*}
    A &= (1 - t)[L_0 + L_1w] - L_0w\\
    B &= L_1 + 2L_2w - 2L_1w
\end{align*}
Considering only what happens for $n \geq 2$ we can collapse the summations
into one summation which is equal to 0, then must be that each term is equal
to 0, so
\begin{align*}
    (n + 1) L_{n + 1} -2 nL_{n} + (n - 1) L_{n-1}
    - (1 - t)L_{n} + L_{n-1} &= 0\\
    (n + 1) L_{n + 1} -2 nL_{n} + nL_{n-1} - L_{n-1}
    - L_n + tL_{n} + L_{n-1} &= 0\\
    (n + 1) L_{n + 1} -(2n + 1 - t)L_{n} + nL_{n-1} &= 0
\end{align*}
The last equation is the relation we are looking for.
\item[(b)] Let us differentiate $\psi$ with respect to $t$
\begin{align*}
    \dv{\psi(t,w)}{t} &= \dv{t}(\frac{1}{1 - w}e^{-tw/(1-w)})\\
    &= -\frac{we^{-tw/(1-w)}}{(1 - w)^2}\\
    &= -\frac{w}{1 - w}\sum_{n=0}^\infty L_n w^n
\end{align*}
From the RHS we have that
\begin{align*}
    \dv{t}(\sum_{n=0}^\infty L_n w^n) &= \sum_{n=1}^\infty L_n' w^{n}
\end{align*}
Where we start from 1 now since $L_0 = 1$ then
\begin{align*}
    -\frac{w}{1 - w}\sum_{n=0}^\infty L_n w^n
    &= \sum_{n=1}^\infty L_n' w^{n}\\
    -w\sum_{n=0}^\infty L_n w^n
    &= (1 - w)\sum_{n=1}^\infty L_n' w^{n}\\
    -\sum_{n=0}^\infty L_n w^{n+1}
    &= \sum_{n=1}^\infty L_n' w^{n}  - \sum_{n=1}^\infty L_n' w^{n+1}\\
    -\sum_{n=1}^\infty L_{n-1} w^{n}
    &= \sum_{n=1}^\infty L_n' w^{n}  - \sum_{n=2}^\infty L_{n-1}' w^{n}\\
    L_0w + \sum_{n=2}^\infty L_{n-1} w^{n}
    &=  - L_1'w + \sum_{n=2}^\infty L_{n-1}' w^{n} - \sum_{n=2}^\infty L_n' w^{n}
\end{align*}
If we only consider the case when $n \geq 2$ then we get the following relation
\begin{align*}
    L_{n-1} &= L_{n-1}' - L_n'
\end{align*}
\end{itemize}
\end{proof}

\cleardoublepage
\begin{proof}{\textbf{13.}}\\
Let us write equation (b) for $n + 1$, we see that
\begin{align*}
    L_n(t) &= L'_n(t) - L'_{n+1}(t)\\
    L'_{n+1}(t) &= L'_n(t) - L_{n}(t)
\end{align*}
Also, taking the derivative of equation (a) with respect to $t$ we get that
\begin{align*}
    (n+1)L'_{n+1}(t) - (2n + 1)L'_n(t) + tL'_n(t) + L_n(t) + nL'_{n-1}(t) &= 0
\end{align*}
Replacing $L'_{n+1}(t)$ we get that
\begin{align*}
    (n+1)(L'_{n}(t) - L_n(t)) - (2n + 1)L'_n(t) + tL'_n(t)
    + L_n(t) + nL'_{n-1}(t) &= 0 \\
    nL'_{n}(t) + L'_{n}(t) - nL_n(t) - L_n(t) - 2nL'_n(t) - L'_n(t) + tL'_n(t)
    + L_n(t) + nL'_{n-1}(t) &= 0 \\
    - nL_n(t) - nL'_n(t) + tL'_n(t) + nL'_{n-1}(t) &= 0
\end{align*}
Now, applying again equation (b) to replace $L'_{n-1}(t)$ we have that
\begin{align*}
    - nL_n(t) - nL'_n(t) + tL'_n(t) + n(L_{n-1}(t) + L'_n(t)) &= 0 \\
    - nL_n(t)  + tL'_n(t) + nL_{n-1}(t) &= 0
\end{align*}
Therefore
\begin{align*}
    tL'_n(t) &= nL_n(t) - nL_{n-1}(t) &  &(c)
\end{align*}
We know that Laguerre's differential equation is 
\begin{align*}
    tL''_{n} + (1 - t)L'_n + nL_n = 0
\end{align*}
Also, differentiating equation (c) with respect to $t$ we have that
\begin{align*}
    tL''_n + L'_n &= nL'_n - nL'_{n-1}
\end{align*}
Replacing this and equation (c), we see that
\begin{align*}
    tL''_{n} + (1 - t)L'_n + nL_n 
    &= nL'_n - nL'_{n-1} - nL_n + nL_{n-1} + nL_n\\
    &= nL'_n - nL'_{n-1} + nL_{n-1}\\
    &= nL'_n - nL'_{n-1} + n(L'_{n-1} - L'_n)\\
    &= 0
\end{align*}
Where we also replaced $L_{n-1}$ with equation (b).\\
We see then that $L_n$ satisfies Laguerre's differential equation.
\end{proof}

\cleardoublepage
\begin{proof}{\textbf{14.}}\\
Let us consider the functions 
\begin{align*}
    e_n(t) = e^{-t/2}L_n(t)
\end{align*}
We want to show they have norm 1 so we need to compute
$\|e_n\| = \sqrt{\inprd{e_n, e_n}}$ as follows
\begin{align*}
    \|e_n\|^2 = \int_0^\infty e^{-t}L_n(t) L_n(t)~dt
\end{align*}
We know that $L_n(t)$ is a polynomial of degree $n$, also, we know from
Problem 15 that if $m < n$ then $\int_0^\infty e^{-t}t^m L_n(t)~dt = 0$
so the only term surviving from $L_n(t)$ in our integral is the term of degree
$n$ i.e.
\begin{align*}
    \|e_n\|^2 &= \int_0^\infty e^{-t}L_n(t) L_n(t)~dt
    = \alpha\int_0^\infty e^{-t}t^n L_n(t)~dt
\end{align*}
Where $\alpha$ is the coefficient of $t^n$. We can determine it using equation
(10c), where we see that $\alpha = (-1)^n/n!$.
Then
\begin{align*}
    \|e_n\|^2 &= \frac{(-1)^n}{n!}\int_0^\infty e^{-t}t^n L_n(t)~dt\\
    &= \frac{(-1)^n}{n!}\int_0^\infty e^{-t}t^n \frac{e^{t}}{n!}f^{(n)}(t)~dt\\
    &= \frac{(-1)^n}{(n!)^2}\int_0^\infty t^nf^{(n)}(t)~dt\\
    &= \frac{(-1)^n}{(n!)^2}\bigg[
    t^nf^{(n-1)}(t)\bigg|_{0}^\infty - n\int_0^\infty t^{n-1}f^{(n-1)}(t)~dt
    \bigg]\\
    &= \frac{n(-1)^{n+1}}{(n!)^2}\int_0^\infty t^{n-1}f^{(n-1)}(t)~dt
\end{align*}
Where we used that $f(t) = t^n e^{-t}$ and we integrated by parts.
So, integrating by parts $n$ times we get that
\begin{align*}
    \|e_n\|^2
    &= \frac{n!(-1)^{2n}}{(n!)^2}\int_0^\infty f^{(0)}(t)~dt
    = \frac{1}{n!}\int_0^\infty t^n e^{-t}~dt
\end{align*}
We apply integration by parts again setting $u = t^n$ and $dv = e^{-t} dt$,
hence
\begin{align*}
    \|e_n\|^2
    &= \frac{1}{n!}\bigg[ -t^ne^{-t}\bigg|_0^\infty
    + n\int_0^\infty t^{n-1} e^{-t}~dt\bigg]
    = \frac{n}{n!} \int_0^\infty t^{n-1} e^{-t}~dt
\end{align*}
Therefore, again applying integration by parts $n$ times we get that
\begin{align*}
    \|e_n\|^2
    = \frac{n!}{n!} \int_0^\infty e^{-t}~dt = 1
\end{align*}
\end{proof}

\cleardoublepage
\begin{proof}{\textbf{15.}}\\
We want to show that the functions 
\begin{align*}
    e_n(t) = e^{-t/2}L_n(t)
\end{align*}
constitute an orthogonal sequence in the space $L^2[0, +\infty)$.\\
First, we want to compute the integral $\int_0^\infty e^{-t}t^m L_n(t)~dt$
where $m < n$.\\
Let us define $f(t) = t^ne^{-t}$ so $L_n(t) = e^tf^{(n)}(t)/n!$, then
\begin{align*}
    \int_0^\infty e^{-t}t^m L_n(t)~dt
    &= \frac{1}{n!}\int_0^\infty t^m f^{(n)}(t)~dt \\
    &= \frac{1}{n!}\bigg[
        t^mf^{(n-1)}(t) \bigg|_0^\infty
        - \int_0^\infty mt^{m-1} f^{(n-1)}(t)~dt
    \bigg]\\
    &= -\frac{m}{n!}\int_0^\infty t^{m-1} f^{(n-1)}(t)~dt
\end{align*}
Where we applied integration by parts and we used that $f^{(n)}(t)$ is of the
form $f^{(n)}(t) = e^{-t} p_n(t)$ where $p_n(t)$ is a polynomial of degree $n$
and when $t \to \infty $ we see that $f^{(n)}(t) \to 0$.\\
Applying integration by parts $m$ times we get that
\begin{align*}
    \int_0^\infty e^{-t}t^m L_n(t)~dt
    &= (-1)^m\frac{m!}{n!}\int_0^\infty f^{(n-m)}(t)~dt\\
    &= (-1)^m(n - m)! \bigg[f^{(n-m-1)}(t)\bigg]_{0}^\infty
\end{align*}
We note that since $n - m - 1 < n$ then the $n-m-1$ derivative of
$f(t) = t^ne^{-t}$ is not going to have yet a constant term so we have
that $f^{(n-m-1)}(t)|_{t=0} = 0$ and we know that $f^{(n-m-1)}(t) \to 0$
when $t \to \infty$, therefore
\begin{align*}
    \int_0^\infty e^{-t}t^m L_n(t)~dt &= 0
\end{align*}
If we take $L_m(t)$ with $m < n$ then $L_m(t)$ is a polynomial of degree
smaller than $n$ and hence
\begin{align*}
    \inprd{e_n(t), e_m(t)} &= \int_0^\infty e^{-t}L_{m}(t) L_n(t)~dt = 0
\end{align*}
Therefore the functions $e_n(t)$ constitute an orthogonal sequence in the space
$L^2[0,+\infty)$.

\end{proof}
\end{document}